% !TeX program = lualatex
% =====================================================================================
%  From Moringa oleifera Radiobiology to Explainable Machine Learning:
%  Cellular Radiosensitivity Prediction and Evidence-Based Compound Prioritisation
%
%  Scientific project report (French Embassy). Final LaTeX edition.
%  Compile with LuaLaTeX + Biber, e.g.  latexmk main.tex  (latexmkrc selects LuaLaTeX),
%  or on Overleaf with the compiler set to LuaLaTeX. See README.md.
% =====================================================================================
\documentclass[11pt,a4paper,oneside,openany,parskip=half-,listof=totoc,
  numbers=noendperiod,captions=tableheading]{scrbook}

\usepackage{embassy-report}
\usepackage{tikz}

\addbibresource{references.bib}
\hypersetup{
  pdftitle={From Moringa oleifera Radiobiology to Explainable Machine Learning:
    Cellular Radiosensitivity Prediction and Evidence-Based Compound Prioritisation},
  pdfsubject={Scientific project report},
  pdfkeywords={radiosensitivity, clonogenic survival, SF2, explainable machine learning,
    radiogenomics, Moringa oleifera, compound prioritisation}}

\begin{document}

% ------------------------------------------------------------------- front matter
\frontmatter
% Cover page and document-information page.
\begin{titlepage}
  \begin{tikzpicture}[remember picture,overlay]
    \fill[reportnavy] (current page.north west) rectangle
      ([yshift=-17mm]current page.north east);
    \node[anchor=west,text=white,font=\sffamily\small]
      at ([xshift=26mm,yshift=-8.5mm]current page.north west)
      {\addfontfeatures{LetterSpace=10}SCIENTIFIC PROJECT REPORT};
    \fill[reportnavy] (current page.south west) rectangle
      ([yshift=5mm]current page.south east);
  \end{tikzpicture}%
  \vspace*{34mm}

  \begin{flushleft}
    {\bfseries\color{reportnavy}\fontsize{26}{32}\selectfont
      From \textit{Moringa oleifera} Radiobiology\\to Explainable Machine Learning\par}
    \vspace{8mm}
    {\color{reportnavy}\rule{0.32\linewidth}{1.2pt}}\par
    \vspace{8mm}
    {\fontsize{15.5}{21}\selectfont\color{reportgrey}
      Cellular Radiosensitivity Prediction and\\
      Evidence-Based Compound Prioritisation\par}
  \end{flushleft}

  \vfill

  \begin{minipage}[b]{0.82\linewidth}
    \small\color{reportgrey}
    A reproducible investigation linking an experimental study of
    \textit{Moringa oleifera} seed and leaf extracts in irradiated cells to an explainable
    benchmark of clonogenic radiosensitivity across 237 independent cancer cell-line
    components, and to a transparent, evidence-based shortlist of plant constituents for
    future experimental validation.
  \end{minipage}

  \vspace{12mm}
  {\color{reportrule}\rule{\linewidth}{0.5pt}}\par\vspace{2mm}
  {\sffamily\small\color{reportgrey}22 October 2026\par}
\end{titlepage}

% ------------------------------------------------------------ document information
\thispagestyle{empty}
\vspace*{\fill}
\begingroup\small
\noindent{\sffamily\bfseries\color{reportnavy}Document information}\par\smallskip
\noindent\begin{tabular}{@{}L{0.27\linewidth}L{0.69\linewidth}@{}}
  \toprule
  Title & From \textit{Moringa oleifera} Radiobiology to Explainable Machine Learning:
    Cellular Radiosensitivity Prediction and Evidence-Based Compound Prioritisation \\
  Document type & Scientific project report, final edition \\
  Date & 22 October 2026 \\
  Primary analysis & Run \path{primary_v1}: five prespecified models, ten repeats of
    grouped five-fold nested cross-validation, 250 of 250 model-fold tasks completed \\
  Further analyses & Runs \path{exploratory_m3_v1}, \path{loto_v1} and \path{scales_v1};
    explainability and triage of \path{primary_v1}; compound evidence and ranking
    \path{compounds_v1} \\
  Not executed & Negative controls (1,000 draws of each null) and the coverage
    simulation for an interval method (Chapter~\ref{ch:implementation}) \\
  Outcome data licence & CLARIFi v1.0, CC BY-NC 4.0; data DOI
    \doilink{10.5281/zenodo.20710748} and the licence must be cited in any reuse \\
  \bottomrule
\end{tabular}

\medskip
\noindent Every number in this report is taken from a saved output of the project
pipeline; the source file of each result table is named in its caption, and the
correspondence with the project's table identifiers is given in
Table~\ref{tab:provenance}. The figures are the vector graphics produced by the
project's figure scripts. No model was refitted for this report.
\par\endgroup
\clearpage

\addchap{Executive summary}\label{ch:summary}

\paragraph{Research motivation.}
An experimental study of crude dichloromethane (DCM) seed and leaf extracts of
\textit{Moringa oleifera} Lam.\ observed that the extracts changed how irradiated cells
proliferated, and that the change differed between an apparently normal rodent fibroblast
line (V79) and two human carcinoma lines (DU145 and MCF-7). A three-line experiment cannot
say whether such differences reflect the cells' \emph{intrinsic} radiation response, nor
whether measurable molecular features explain it. This project asked that broader question
on a defensible endpoint and with public data, and then returned to the plant to ask which
of its documented constituents deserve experimental testing.

\paragraph{Computational question.}
Does baseline gene expression, summarised as 50 Hallmark pathway scores, predict
experimentally measured clonogenic radiosensitivity across cancer cell lines better than
prespecified simple baselines, once related cell lines are kept together and tissue of
origin is already taken into account?

\paragraph{Data and design.}
We assembled 241 eligible cell-line records, forming 237 independent biological components,
by linking literature-derived clonogenic surviving fractions at 2~Gy (CLARIFi mSF2) to
baseline expression (DepMap 24Q4). Cell-line identity and relatedness were resolved through
Cellosaurus, and pathway features were computed from the MSigDB Hallmark gene sets. Five
prespecified models, from a global mean to tissue, proliferation and expression together,
were compared under ten repeats of five-fold cross-validation grouped by biological
component and stratified by tissue, with every tuning and scaling step confined to training
folds. All 250 model-fold tasks completed.

\paragraph{Principal findings.}
Adding expression features to a tissue-only baseline reduced the mean absolute error (MAE)
by \textbf{0.00383 mSF2 units, 2.85\% of the baseline error}, and the richer model had the
lower error in \textbf{nine of ten} cross-validation repeats. Out-of-sample $R^2$ stayed
close to zero for every model ($-0.018$ to $+0.013$). Expression therefore carries a small
amount of usable information beyond tissue of origin, while overall predictive power is
weak --- a result that agrees with the strongest published benchmark in the field.
Superiority was not formally tested: a protocol deviation, recorded before any result was
seen, deferred the coverage simulation needed to validate an interval method, and no
negative-control draws were run. The improvement is reported as a descriptive estimate.

\paragraph{Robustness and generalisation.}
Three further analyses of the same frozen cohort tested how durable this small advantage
is. A random forest (M3) did \textbf{worse} than the linear offset model (gain over the
tissue baseline 0.00275 against 0.00383). On transformed outcome scales the advantage kept
its direction but almost vanished on the neglog scale (0.00023 in neglog units) and was
larger in absolute terms on the logit scale (0.00980 in logit units), so its size depends on
how the outcome is expressed. When one tissue at a time was held out, an expression-only model
reached out-of-sample $R^2$ of \textbf{+0.031} against a global mean, the only clearly
positive $R^2$ in the project; it nevertheless lost to the global mean when Pancreas
($-0.0142$) or Esophagus/Stomach ($-0.0030$) was held out, so transfer to an unseen tissue
is conditional on the tissue.

\paragraph{Explainability.}
Held-out permutation importance ranked the 50 Hallmark scores. Seven reached the
prespecified exploratory tier~C, led by the p53 pathway; the higher tiers were unreachable
by design, because they require a formally demonstrated prediction gain and at least 1,000
negative-control draws. Importance magnitudes are as small as the predictive gain itself,
rankings are only moderately stable across fits, and the DNA repair, G2M checkpoint and
hypoxia sets that radiobiology would lead one to expect were not prioritised.

\paragraph{Compound prioritisation.}
\textit{M.~oleifera} seed and leaf constituents were curated from primary analytical
studies, with the extraction solvent of every report recorded, and linked to
radiation-response biology through a four-dimension score of 0--8 points fixed before any
score was computed. Eighteen compounds were ranked and five reached tier~A. The curation
exposed the central mismatch of the project: \textbf{niazimicin is the only constituent
with a structurally confirmed identification in a dichloromethane fraction of both seed
and leaf, yet no study has tested the purified compound in any radiation-response process;
conversely, the compounds with the strongest radiation evidence have never been found in a
dichloromethane extract.} Two prespecified safeguards visibly bound: chlorogenic acid scored
7 of 8 but is held at tier~C because one study reports that it protects \emph{tumour} cells
from radiation, and benzyl isothiocyanate is excluded despite strong radiation evidence
because it is a root, not a seed or leaf, constituent. A tier summarises the state of the
published evidence; it never asserts radioprotection.

\paragraph{Significance.}
The project delivers a reproducible, auditable pipeline, an honest benchmark whose
negative-leaning result agrees with the best prior work, and a traceable route from the
original radiobiology to a short list of experiments that could settle the question. Its
value lies in the rigour of its methods and the transparency of its evidence rather than in
a strong performance claim.

\begin{keybox}[title={Recommended next steps.}]
Characterise the original DCM extracts analytically (niazimicin and moringin first);
complete the deferred coverage simulation and the negative controls; and test
$\alpha$-tocopherol and $\beta$-sitosterol, purified and against whole extract, on
clonogenic survival in matched normal and tumour cells over a dose range.
\end{keybox}

\tableofcontents
\listoffigures
\listoftables

% -------------------------------------------------------------------- main matter
\mainmatter
\chapter{Introduction and scientific background}\label{ch:introduction}

\section{Cellular responses to ionising radiation}

Ionising radiation kills cells mainly by damaging DNA, and the DNA double-strand break is
the lesion most closely tied to loss of reproductive capacity. The reference measure of how
much radiation a cell population survives is the \emph{clonogenic assay}, which counts the
cells that retain unlimited proliferative capacity and form colonies after irradiation
\parencite{puck1956}. For a dose $D$, the surviving fraction is the plating efficiency of
irradiated cells relative to that of unirradiated controls,
\begin{equation}
  \mathrm{SF}(D) \;=\; \frac{N_{\text{colonies}}(D)\,/\,N_{\text{seeded}}(D)}
                             {N_{\text{colonies}}(0)\,/\,N_{\text{seeded}}(0)} .
  \label{eq:sf}
\end{equation}
The surviving fraction at 2~Gy (SF2) is the most widely used single summary, because 2~Gy
is the conventional daily fraction in radiotherapy and because SF2 has been associated with
clinical outcome in some settings \parencite{fertil1981,west1993}.

Cell lines differ substantially and reproducibly in intrinsic radiosensitivity, and these
differences are only partly explained by tissue of origin \parencite{yard2016}. The
determinants include DNA damage response and repair capacity, cell-cycle checkpoint
control, proliferation rate, apoptotic propensity, oxygenation and redox state. Each of
these processes leaves a transcriptional footprint, which is why prediction from gene
expression is plausible in principle.

\section{The \textit{Moringa oleifera} experimental motivation}\label{sec:moringa-study}

The motivating study is an unpublished manuscript, \enquote{Dichloromethane seed and leaf
extracts of \textit{Moringa oleifera} (Lam) enhance proliferation in irradiated rodent
fibroblasts \textit{in vitro}}. The scanned copy available to the project is incomplete:
printed pages 1, 2, 4, 5, 6, 13, 14 and 15 are present, while page~3 and pages 7--12 are
missing. The plant material and extraction methods, the end of the discussion, the
conclusions and the reference list therefore cannot be read; authors, venue, date and DOI
are not shown, and no published version was located. No missing method, result or
conclusion has been reconstructed in this report.

The available pages document the following design. Crude DCM seed and leaf extracts were
reconstituted in phosphate-buffered saline, sterile-filtered and applied at
0--100~µg/mL to V79 (described as apparently normal Chinese hamster lung fibroblasts),
DU145 (human prostate cancer) and MCF-7 (human mammary adenocarcinoma). For each cell line
and extract, one plate was left unirradiated and a second received 2~Gy of X-rays at
1~Gy/min. Adherent biomass was measured 72~h later by crystal-violet optical density at
590~nm, with dose-matched untreated controls subtracted. The authors integrated the change
in optical density over extract concentration, giving an area $A_D$ at each dose $D$, and
reported a \emph{proliferation enhancement factor}
\begin{equation}
  \mathrm{PEF} \;=\; \frac{A_{2\,\mathrm{Gy}}}{A_{0\,\mathrm{Gy}}} .
  \label{eq:pef}
\end{equation}
For the seed extract, PEF exceeded 1 in all three lines (V79 1.25, DU145 1.12, MCF-7 1.16).
For the leaf extract, PEF was 1.26 in V79, 0.90 in DU145 and 0.99 in MCF-7.

Four cautions follow directly from the design, and they constrain everything that follows.
\begin{enumerate}
  \item \textbf{PEF is not SF2.} PEF is a ratio of integrated extract-versus-control
    differences in 72-hour biomass measured with the extract present; SF2 is a surviving
    fraction of untreated cells measured over one to three weeks of colony growth. The two
    are never equated, pooled or plotted on one axis in this project.
  \item \textbf{\enquote{Enhancement} denotes two different quantities in the manuscript.}
    By PEF the seed extract enhanced all three lines; by growth relative to untreated
    control only V79 increased, while the two carcinoma lines showed no change or reduced
    growth. The discussion's statement about the fibroblasts refers to the second quantity,
    the results' statement about all lines to the first.
  \item \textbf{The normal-versus-malignant comparison is confounded.} The only apparently
    normal line is also the only non-human line, so species, tissue of origin, phenotype
    and malignancy vary together. Only the leaf extract shows a normal/malignant difference
    in PEF.
  \item \textbf{The design cannot support a selectivity claim.} There is a single non-zero
    dose and no dose--response; the replicates are technical wells rather than independent
    biological repeats; dose is confounded with plate; and no vehicle or cell-free
    interference control is described. The available pages contain no clonogenic,
    cytotoxicity, apoptosis, DNA-damage or migration assay, no normal human cell line and
    no mechanistic experiment.
\end{enumerate}
The \textit{Moringa} measurements are therefore used as \emph{biological motivation only}.
They are never used as training, calibration or validation data, and no claim of selective
phytochemical radioprotection is made anywhere in this report.

\section{Why cellular radiosensitivity is hard to predict}

Published attempts to predict cell-line radiosensitivity from expression have performed
poorly under honest evaluation. The most directly relevant prior work, RadSigBench,
re-fitted seven published radiosensitivity signatures and benchmarked them against
resampled random gene sets, reporting a mean out-of-sample $R^2$ of 0.01 on CCLE data and
0.00 on NCI-60, with the published signatures performing no better than random gene sets
\parencite{oconnor2022}. This is the single most important prior for the present work: it
sets the expectation that the signal is weak, so a near-zero $R^2$ here is a replication
rather than a surprise.

Three difficulties recur. The first is \emph{sample size}: panels with reliable clonogenic
data number in the hundreds, against tens of thousands of candidate molecular features. The
second is \emph{label noise}: literature-aggregated SF2 values combine laboratories, dose
rates and protocols, and an audit of 1,672 radiobiology papers found that 30.5\% did not
report replicates and 32.3\% did not report dose rate \parencite{oike2020}. The third is
\emph{leakage through relatedness}: isogenic derivatives, sublines and cross-contaminated
lines share molecular profiles, so a split that separates records rather than biological
entities inflates apparent performance.

\section{The case for public radiogenomic data}

A three-line extract experiment cannot address intrinsic radiosensitivity across cell
types, but public resources can. CLARIFi aggregates literature clonogenic radiosensitivity
\parencite{clarifi2024}; DepMap provides harmonised baseline expression for the same lines
\parencite{depmap24q4}; Cellosaurus provides authoritative identity, problem flags and
relatedness \parencite{bairoch2018}; and the MSigDB Hallmark collection provides 50 compact,
interpretable pathway summaries \parencite{liberzon2015}. Together they allow a properly
grouped, prespecified benchmark on a defensible endpoint.

\section{The value and the limits of interpretable machine learning}

Interpretable models matter here for two reasons: a pathway-level summary is biologically
meaningful, and feature attributions can be examined for stability. Attribution, however,
has hard limits. Permutation importance describes a \emph{fitted model}, not biology;
correlated features share or exchange importance; and the size of an importance is not a
molecular effect size. Attribution is only as trustworthy as the prediction it explains,
which is why every pathway statement in this report is capped at an exploratory tier while
the prediction gain remains small and unvalidated.

\section{Rationale for prioritising candidate compounds}

The original study raises a question that the machine-learning benchmark alone cannot
answer: if cells differ in radiation response, and some of that difference is reflected in
pathway activity, which \textit{Moringa} constituents are worth testing? The answer has to be
built compound by compound from independently published chemistry and biology, with the
extraction solvent tracked explicitly, because the extracts used in the original study were
never chemically characterised.

\chapter{Aim and objectives}\label{ch:aim}

\paragraph{Aim.}
To build a reproducible, auditable research pipeline that connects the original
\textit{Moringa oleifera} irradiation study to a transparent, explainable benchmark of
cellular radiosensitivity prediction on public radiogenomic data, and to use that benchmark,
together with curated published evidence, to produce a traceable and explicitly exploratory
shortlist of candidate constituents and a feasible strategy for their experimental
validation.

Figure~\ref{fig:workflow} shows the six stages of the project and the order in which they
were carried out. Stages 1--5 are complete; stage 6 is a proposal.

\begin{figure}[tbp]
  \centering
  \includegraphics[width=0.86\linewidth,height=\ReportFigureMaxHeight,keepaspectratio]{fig_01_project_workflow}
  \caption[Structure of the project]{Structure of the project, from the original
    \textit{Moringa oleifera} radiobiology investigation to the prioritisation of candidate
    constituents for future experimental work. Box fills distinguish data sources,
    computational steps, literature curation and proposed future work; they do not indicate
    the strength of the evidence at each stage. The \textit{Moringa} measurements are
    biological motivation only and are never used as training, calibration or validation
    data, because crystal-violet proliferation under extract treatment and clonogenic
    survival at 2~Gy are different endpoints. Stage 6 is a proposal, not completed work.}
  \label{fig:workflow}
\end{figure}

Table~\ref{tab:objectives} lists the ten objectives and their status. Nine are complete.
The validation of the primary comparison by a coverage simulation and negative controls
(objective 9) was deferred by a protocol deviation recorded before any result existed
(Section~\ref{sec:deviation}); three other robustness analyses were executed in its place
(Chapter~\ref{ch:robustness}).

\begin{table}[tbp]
  \caption{Project objectives and their status.}\label{tab:objectives}
  \small
  \begin{tabular}{@{}r L{0.60\linewidth} L{0.30\linewidth}@{}}
    \toprule
    \thead{\#} & \thead{Objective} & \thead{Status} \\
    \midrule
    1 & Extract and document the original manuscript's methods and results, marking every
        unreadable or absent item & Completed (Section~\ref{sec:moringa-study}) \\
    2 & Assemble a provenance-tracked dataset linking clonogenic SF2 to baseline
        expression, with cell-line identity and relatedness resolved
      & Completed; 241 records, 237 components \\
    3 & Prespecify and execute a grouped, nested cross-validated comparison of five models
        on a single named endpoint & Completed; 250/250 model-fold tasks \\
    4 & Quantify the predictive gain of expression beyond tissue of origin and report it
        with its limitations & Completed; descriptive estimate \\
    5 & Rank Hallmark gene-set scores by held-out permutation importance under
        prespecified evidence gates & Completed; capped at tier~C \\
    6 & Triage the cell lines whose measured radiosensitivity the model reproduces least
        well, separating technical from biological explanations & Completed \\
    7 & Curate \textit{M.~oleifera} seed and leaf constituents and their radiation-response
        evidence from primary sources, tracking the extraction solvent
      & Completed (Tables~\ref{tab:c1} and \ref{tab:c2}) \\
    8 & Apply a transparent 0--8 point prioritisation fixed before computation and produce
        the ranked shortlist & Completed; 18 ranked, 5 in tier~A \\
    9 & Validate the comparison with a coverage simulation and negative controls
      & Deferred by recorded protocol deviation \\
    10 & Specify a feasible experimental validation strategy
      & Completed as a proposal (Chapter~\ref{ch:recommendations}) \\
    \bottomrule
  \end{tabular}
\end{table}

\chapter{Data sources and resources}\label{ch:data}

The resources used by the project are listed in Table~\ref{tab:resources}. Every raw file
carries a manifest recording its source, DOI, URL, access date, retrieval method, SHA-256
checksum, size, column meanings, units, endpoint definition and transformations; the data
loaders verify the checksum and stop on any mismatch. The outcome file, the expression
matrix, the Hallmark gene-set file and the quality-control outputs were all verified
against their pinned checksums before the analysis run.

\begin{table}[tbp]
  \caption{Data resources, releases, identifiers and licences.}\label{tab:resources}
  \footnotesize
  \begin{tabular}{@{}L{0.17\linewidth}L{0.13\linewidth}L{0.27\linewidth}L{0.15\linewidth}L{0.20\linewidth}@{}}
    \toprule
    \thead{Resource} & \thead{Release} & \thead{Identifier} & \thead{Licence}
      & \thead{Role in the project} \\
    \midrule
    Original \textit{Moringa} manuscript & Unpublished scan, 8 pages
      & SHA-256 \path{7f644f55…e870fb9} & Redistribution not authorised
      & Biological motivation only \\
    CLARIFi radiosensitivity & v1.0
      & Data \doilink{10.5281/zenodo.20710748}; paper \doilink{10.1111/cas.16334}
      & CC BY-NC 4.0 & Outcome: mSF2 \\
    DepMap & 24Q4 Public & \doilink{10.25452/figshare.plus.27993248.v1} & CC BY 4.0
      & Baseline expression; cell-line metadata \\
    Cellosaurus & r56.0 & \doilink{10.26037/yareta:lbz6rhlaavbldmavv4k5jdaphe}
      & CC BY 4.0 & Identity, problem flags, relatedness \\
    MSigDB Hallmark & 2025.1.Hs & Paper \doilink{10.1016/j.cels.2015.12.004}
      & MSigDB terms & 50 pathway features \\
    PubChem, Europe PMC, Crossref & Accessed 22 Oct 2026 & \na & Public interfaces
      & Compound identity; literature verification \\
    \bottomrule
  \end{tabular}
\end{table}

\paragraph{Licensing and provenance caveats.}
\begin{itemize}
  \item \textbf{CLARIFi is licensed CC BY-NC 4.0} (non-commercial), as confirmed from both
    the Crossref licence record and the Europe PMC full text. The raw file is kept outside
    version control, the data DOI and the licence must be cited in any publication, and the
    non-commercial terms constrain downstream reuse.
  \item The \textbf{Hallmark gene-set file} was downloaded manually, so no acquisition log
    or published checksum exists for it. Its content was verified locally (48,689 bytes,
    50 unique \texttt{HALLMARK\_} sets, standard layout, SHA-256 \path{f22066af…cbbdd8e9}),
    but its identity with the official release has not yet been confirmed against the
    official download address.
  \item The \textbf{original manuscript may not be redistributed}; only its checksum and
    the page-by-page extraction record are kept.
  \item Compound curation used public databases and published literature only; no licensed
    chemical database was queried.
\end{itemize}

\chapter{Methodology}\label{ch:methods}

\section{The original experimental record}

The \textit{Moringa} manuscript was read page by page, and every extracted item was labelled
as a measured result, an authors' interpretation, this project's commentary, or not stated;
nothing absent was reconstructed. A radiobiology review checked every extracted item against
the page images. Internal inconsistencies were recorded as printed rather than silently
corrected: the extract concentration appears as 25--50~µg/mL in the abstract, 30--50~mg/mL
in the results and above 20~µg/mL in the discussion; V79 is described variously as hamster,
murine and rodent; and the tabulated integrated areas are all positive although the plotted
optical-density changes are mostly negative for the two carcinoma lines, which cannot be
reconciled without the missing pages or the raw data.

\section{Data acquisition, harmonisation and eligibility}

Outcome, expression, identity and pathway resources were acquired under recorded
provenance, parsed into separate per-endpoint tables and harmonised on DepMap identifiers.
Download code was written only after each source had been verified; raw files are
immutable, no value was invented or silently imputed, and missingness is documented.

Cell lines were matched by stable identifiers, never by name. Cellosaurus supplied
accessions, synonyms, problem flags, and the hierarchy and origin links used to detect
relatedness. A record was eligible if it had a usable clonogenic SF2 value, a matched
baseline expression profile and a resolved identity, which gave \textbf{241 eligible
records} from 20 tissues.

\section{Endpoint and features}\label{sec:endpoint}

The single primary endpoint is \textbf{CLARIFi mSF2}, the median of the literature-reported
clonogenic surviving fractions at 2~Gy for a cell line, analysed on its raw scale.
Predictions are clipped to $[0,1]$ by a prespecified rule; outcomes are never clipped.
Crystal-violet proliferation, cell viability, ATP-based integral survival and differential
expression are kept strictly separate from clonogenic SF2 throughout the project.

Three kinds of feature were used: tissue of origin (OncotreeLineage, one-hot encoded over
fixed categories), an optional proliferation covariate, and 50 Hallmark gene-set scores
computed from baseline expression (DepMap $\log_2(\mathrm{TPM}+1)$ values of protein-coding
genes). Within each training fold, every gene $g$ is standardised with the mean $\mu_g$ and
standard deviation $\sigma_g$ of the training rows only, and the score of set $S$ for cell
line $i$ is the mean standardised expression of its member genes present in the matrix,
\begin{equation}
  s_{iS} \;=\; \frac{1}{|S^{\ast}|}\sum_{g\in S^{\ast}} \frac{x_{ig}-\mu_g}{\sigma_g},
  \qquad S^{\ast} = S \cap \{\text{measured genes}\}.
  \label{eq:score}
\end{equation}
The proliferation covariate is the combined score of the E2F-targets and G2M-checkpoint
sets.

\section{Grouping and cross-validation}\label{sec:cv}

Records of cell lines related by identity (isogenic derivatives, sublines, documented
relatedness) were merged into \textbf{237 independent biological components}; seven
components contain more than one record. Every split is grouped by component, so related
lines can never fall on both sides of a train/test boundary. Splits are stratified by
tissue: the nine tissues with at least ten components each form their own stratum and the
remaining eleven are pooled.

Validation used \textbf{ten repeats of five-fold cross-validation} with an inner five-fold
loop for hyperparameter selection. All preprocessing, scaling, gene-set scoring and tuning
took place \textbf{inside training folds only}. Each record $i$ belonging to component
$c(i)$ carries the weight
\begin{equation}
  w_i \;=\; \frac{1}{n_{c(i)}},
  \label{eq:weights}
\end{equation}
where $n_c$ is the number of records of component $c$, in every fit and every inner score,
so that a component with several records counts once. The base random seed (20261021) and
all grids are recorded in the experiment configuration. A dedicated leakage audit covered
identifier handling, component construction and split construction.

\section{Models}\label{sec:models}

Five models were prespecified, from simple to rich (Table~\ref{tab:models}), and simple
scikit-learn estimators were used before any more complex architecture. M1 is a ridge
regression on the fixed tissue categories with an unpenalised intercept, which shrinks each
tissue mean towards the overall mean. M2 is an \emph{offset model}: within each training
fold it first fits M1 and then an elastic net to the M1 residuals using the standardised
gene-set scores $\tilde{\mathbf{s}}_i$,
\begin{equation}
  \hat{y}_i^{\,\mathrm{M2}} \;=\; \hat{y}_i^{\,\mathrm{M1}} + \beta_0 +
     \tilde{\mathbf{s}}_i^{\top}\boldsymbol{\beta},
  \qquad
  \hat{\boldsymbol{\beta}} = \arg\min_{\beta_0,\boldsymbol{\beta}}
     \frac{\sum_i w_i\bigl(r_i-\beta_0-\tilde{\mathbf{s}}_i^{\top}\boldsymbol{\beta}\bigr)^2}
          {2\sum_i w_i}
     + \alpha\rho\,\lVert\boldsymbol{\beta}\rVert_1
     + \frac{\alpha(1-\rho)}{2}\lVert\boldsymbol{\beta}\rVert_2^2 ,
  \label{eq:offset}
\end{equation}
with $r_i = y_i-\hat{y}_i^{\,\mathrm{M1}}$. Because M2 contains M1 as an offset, the
comparison of the two asks specifically what expression adds beyond tissue. M1P and M2P are
the corresponding models with the proliferation covariate added to the baseline; M2P uses
the remaining 48 Hallmark scores in its residual model. The inner loop searched
$\alpha\in\{0.1, 0.01, 0.001\}$, $\rho\in\{0.9, 0.5, 0.1\}$ and the baseline ridge penalty
$\in\{100, 10, 1, 0.1\}$, 36 combinations for M2.

\begin{table}[tbp]
  \caption{The five prespecified models of the primary benchmark.}\label{tab:models}
  \small
  \begin{tabular}{@{}l L{0.50\linewidth} L{0.30\linewidth}@{}}
    \toprule
    \thead{Model} & \thead{Features} & \thead{Role} \\
    \midrule
    M0 & None (training-set mean) & Floor \\
    M1 & Tissue of origin & Primary comparator \\
    M1P & Tissue and proliferation score & Sensitivity comparator \\
    \textbf{M2} & \textbf{Tissue and 50 Hallmark expression scores} & \textbf{Primary model} \\
    M2P & Tissue, proliferation and the other 48 Hallmark scores & Sensitivity model \\
    \bottomrule
  \end{tabular}
\end{table}

Two further models were configured for the exploratory layer
(Chapter~\ref{ch:robustness}): M2X, an elastic net on the 50 scores with the global mean as
offset and no tissue information, used only for the leave-one-tissue-out analysis; and M3,
a random forest on one-hot tissue and the 50 scores. M3 was added after the primary results
had been seen and therefore cannot replace M2 as the primary benchmark, whatever its score.

\section{Performance measures and the primary estimand}\label{sec:metrics}

For a repeat $r$ with out-of-fold predictions $\hat{y}_{ir}$, the component-weighted error
measures are
\begin{align}
  \mathrm{MAE}_r &= \frac{\sum_i w_i\,\lvert y_i-\hat{y}_{ir}\rvert}{\sum_i w_i}, &
  \mathrm{RMSE}_r &= \Biggl(\frac{\sum_i w_i\,(y_i-\hat{y}_{ir})^2}{\sum_i w_i}\Biggr)^{1/2},
  \label{eq:mae}\\
  R^2_{\mathrm{oos},r} &= 1-\frac{\sum_i w_i\,(y_i-\hat{y}_{ir})^2}
                                  {\sum_i w_i\,(y_i-\bar{y}_w)^2}, &
  \bar{y}_w &= \frac{\sum_i w_i y_i}{\sum_i w_i},
  \label{eq:r2}
\end{align}
and each reported value is the mean over the $R=10$ repeats. $R^2_{\mathrm{oos}}=0$
therefore corresponds to predicting the cohort mean, and negative values are worse than
that. Calibration slope and intercept come from a weighted regression of observed on
predicted values.

The primary estimand is the paired, component-weighted difference in mean absolute error
between M1 and M2. With the record-level difference
$d_{ir}=\lvert y_i-\hat{y}^{\,\mathrm{M1}}_{ir}\rvert-\lvert y_i-\hat{y}^{\,\mathrm{M2}}_{ir}\rvert$,
the component means $D_{cr}=n_c^{-1}\sum_{i\in c} d_{ir}$ and the $C=237$ components,
\begin{equation}
  \Delta\mathrm{MAE} \;=\; \frac{1}{R}\sum_{r=1}^{R}\,\frac{1}{C}\sum_{c=1}^{C} D_{cr},
  \label{eq:dmae}
\end{equation}
so a positive value favours the richer model. MAE, RMSE and $R^2_{\mathrm{oos}}$ are
reported for every model, and no metric was chosen after the results were seen.

\section{Explainability}\label{sec:explain-methods}

For each of the 50 real-data fits of M2 (ten repeats by five folds), grouped permutation
importance was computed on the \emph{outer-test rows only}, using the saved model, with 20
permutations per feature; nothing was refitted. For gene set $S$ in fit $f$,
\begin{equation}
  I_{S,f} \;=\; \frac{1}{20}\sum_{p=1}^{20}
     \Bigl[\mathrm{MAE}_f\bigl(\text{score } S \text{ permuted across components, draw } p\bigr)
           -\mathrm{MAE}_f\Bigr],
  \label{eq:importance}
\end{equation}
the increase in component-weighted held-out error, in mSF2 units, when the information in
that score is destroyed. Coefficient sign consistency across fits and the Spearman
correlation between gene-set scores were also computed, so that shared signal between
correlated sets is visible rather than hidden.

\section{Pathway prioritisation}\label{sec:pathway-methods}

Gene sets were assigned to tiers by rules frozen before the run. \emph{Tier~C} requires that
importance exceeded zero in at least half the fits and that the share of fits with a zero
coefficient was at most 0.50. \emph{Tier~B} additionally requires sign consistency and a
positive-importance share of at least 0.75 and significance against the permuted-outcome
null; \emph{tier~A} requires both at least 0.90 and significance against both null models.
\textbf{Both higher tiers also require a formally demonstrated prediction gain and at least
1,000 completed draws of each null.} Neither exists in this run, so tiers A and B were
unreachable by construction, and the rules were not altered to produce a more impressive
ranking. Scores correlated at $\lvert\rho\rvert\ge 0.80$ are flagged as sharing signal.

\section{Cell-line triage}\label{sec:triage-methods}

Each measured line was characterised by the absolute error of its repeat-averaged
out-of-fold prediction under M1 and M2, the spread of its prediction across repeats, and
the disagreement between the two models, each expressed as a cohort percentile (the share
of lines with a strictly smaller value, plus one, divided by the number of lines). A line
reaches \emph{tier~A} when its error is at or above the 95th percentile under both models,
\emph{tier~B} at or above the 90th percentile under both, and \emph{tier~C} at or above the
90th percentile under either. Technical flags --- stranded RNA profile, related lines in the
cohort, a Cellosaurus problem flag on a relative, a small tissue group, unstable prediction
and strong model disagreement --- are recorded separately from the tier, so that a large
residual caused by a data or identity problem is never presented as a biological
discrepancy.

\section{Curation of \textit{Moringa} constituent evidence}\label{sec:curation-methods}

Constituents and their biology were curated by literature search and reading only, into two
validated evidence files. Only \emph{primary} studies count, so a review citing an earlier
study is never independent confirmation. Only studies of the \textbf{purified compound} can
score; a study that dosed an extract is recorded with its score contribution disabled. Only
\emph{ionising} radiation qualifies as radiation evidence, which excludes the large UV
photoprotection literature. Nanoparticle and metal-conjugate formulations are excluded
because the carrier confounds the compound, and chemically distinct derivatives are never
credited to the parent compound. The \textbf{extraction solvent of every report is recorded
verbatim}, and each compound is classified as documented in a dichloromethane extract,
documented in seed or leaf with another solvent, or indirect and unconfirmed. Every DOI was
verified against Europe PMC or Crossref and every identifier against PubChem, with the
verification route and reading depth stored per record.

\section{Candidate scoring and tiers}\label{sec:scoring}

Each compound receives four integer scores, fixed before any score was computed:
\begin{equation}
  T \;=\; \underbrace{D_1}_{0\text{--}2}+\underbrace{D_2}_{0\text{--}3}
         +\underbrace{D_3}_{0\text{--}2}+\underbrace{D_4}_{0\text{--}1}\;\in\{0,\dots,8\},
  \label{eq:priority}
\end{equation}
where $D_1$ is the strength of verified \textit{Moringa} constituent evidence (2 if at least
one primary seed or leaf report has structural confirmation by isolation and NMR or by a
reference standard; 1 if identification rests on a spectral-library match or is not stated;
0 if no primary seed or leaf report exists); $D_2$ the relevance to radiation-response
biology (3 for ionising-radiation evidence from at least two independent research groups,
2 from one group, 1 if only a relevant process was modulated without radiation); $D_3$
independent source support (2 for at least three distinct research groups across both
evidence files, 1 for two); and $D_4$ feasibility of follow-up (1 if the identifier is
verified and at least one scored study reports a usable concentration or dose). Research
groups are counted once across both evidence files, so repeated work from one laboratory is
one source, and \textbf{chemical-class analogy earns no points}.

The tiers are A ($T\ge 6$, $D_1\ge 1$, $D_2\ge 2$), B ($T\ge 4$, $D_1\ge 1$, $D_2\ge 1$) and C
($T\ge 2$, $D_1\ge 1$); compounds below these thresholds are not prioritised. Two safeguards
override the score. A compound with no primary \textit{Moringa} identification ($D_1=0$) is
\emph{not assessable} and never enters the shortlist. Evidence that a compound protects
\emph{tumour} cells from radiation, or is itself genotoxic, is a liability that is displayed
and caps the tier at C. Extraction-specific evidence is reported beside the score and never
folded into it.

Model-associated pathways, literature-supported processes and their overlap are kept apart.
An overlap between the genes a published experiment measured and the members of a
model-ranked gene set marks a hypothesis to be tested; it is not evidence that the compound
modifies the gene set, the pathway or clonogenic SF2.

\section{Reproducibility and quality assurance}\label{sec:repro}

The codebase uses Python~3.11 or later with a \path{src/} layout, type hints, small pure
functions and YAML configuration with recorded seeds and no machine-specific paths. Every
analytical output carries a run record with the git commit and a flag for uncommitted
changes, configuration and file checksums, package versions and UTC timestamps, and
analysis runs refuse to start from an uncommitted working tree. Result folders and figures
are write-once, so a published number cannot be silently overwritten. Specialist reviews
covering radiobiology, radiogenomic data curation, machine-learning methodology, statistics,
software architecture, reproducibility, figure design and literature were recorded at each
decision point in the project's decision log, together with any unresolved disagreement.

\chapter{Machine-learning experiments and results}\label{ch:results}

\section{The analysed cohort}\label{sec:cohort}

The analysis cohort comprises \textbf{241 eligible records forming 237 independent
biological components} across 20 tissues (Table~\ref{tab:composition};
Figure~\ref{fig:cohort}a). Nine tissues have at least ten components and form their own
cross-validation strata: Lung (46), CNS/Brain (23), Esophagus/Stomach (23), Bowel (21), Head
and Neck (18), Breast (17), Pancreas (16), Bladder/Urinary Tract (11) and Ovary/Fallopian
Tube (10). The remaining eleven tissues, from Skin (9) to Biliary Tract (1), are pooled.
Seven components contain more than one record, and seven expression profiles are flagged
as stranded RNA preparations.

\begin{table}[tbp]
  \caption[Dataset composition by tissue]{Dataset composition by tissue of origin
    (OncotreeLineage). Components group related cell lines; tissues with fewer than ten
    components are pooled into one cross-validation stratum. Source:
    \path{results/tables/primary_v1/publication_v1/table_1_dataset_composition.csv}.}
  \label{tab:composition}
  \small
  \begin{tabular}{@{}L{0.30\linewidth}S[table-format=2.0]S[table-format=3.0]S[table-format=1.0]S[table-format=1.0]L{0.20\linewidth}@{}}
    \toprule
    \thead{Tissue} & {\thead{Records}} & {\thead{Components}}
      & {\thead{Multi-record}} & {\thead{Stranded RNA}} & \thead{Stratification} \\
    \midrule
    Lung & 46 & 46 & 1 & 1 & own stratum \\
    CNS/Brain & 23 & 23 & 1 & 2 & own stratum \\
    Esophagus/Stomach & 23 & 23 & 0 & 0 & own stratum \\
    Bowel & 23 & 21 & 2 & 1 & own stratum \\
    Head and Neck & 18 & 18 & 0 & 0 & own stratum \\
    Breast & 17 & 17 & 0 & 0 & own stratum \\
    Pancreas & 16 & 16 & 0 & 1 & own stratum \\
    Bladder/Urinary Tract & 11 & 11 & 0 & 0 & own stratum \\
    Ovary/Fallopian Tube & 10 & 10 & 0 & 0 & own stratum \\
    \addlinespace
    Skin & 9 & 9 & 0 & 0 & pooled \\
    Bone & 9 & 8 & 1 & 1 & pooled \\
    Uterus & 9 & 8 & 1 & 0 & pooled \\
    Pleura & 6 & 6 & 0 & 0 & pooled \\
    Prostate & 5 & 5 & 0 & 1 & pooled \\
    Kidney & 3 & 3 & 0 & 0 & pooled \\
    Liver & 3 & 3 & 0 & 0 & pooled \\
    Peripheral Nervous System & 3 & 3 & 1 & 0 & pooled \\
    Soft Tissue & 3 & 3 & 0 & 0 & pooled \\
    Thyroid & 3 & 3 & 0 & 0 & pooled \\
    Biliary Tract & 1 & 1 & 0 & 0 & pooled \\
    \midrule
    \textbf{All} & \bfseries 241 & \bfseries 237 & \bfseries 7 & \bfseries 7 & \\
    \bottomrule
  \end{tabular}
\end{table}

Observed mSF2 spans almost the whole possible interval, and tissue medians differ visibly
(Figure~\ref{fig:cohort}b), which is why a tissue-only model is the appropriate baseline
against which to judge expression features. Two of the three cell lines of the motivating
study are present: MCF-7 (observed mSF2 0.52) and DU145 (0.71). They are shown as context
only. V79 is a Chinese hamster line and has no counterpart in this human cohort.

\begin{figure}[tbp]
  \centering
  \includegraphics[width=\linewidth,height=\ReportFigureMaxHeight,keepaspectratio]{fig_02_dataset_composition}
  \caption[Composition and radiosensitivity of the analysed cohort]{Composition and
    radiosensitivity of the analysed cohort (241 eligible records, 237 independent
    biological components). (a)~Independent components per tissue (OncotreeLineage); the
    dashed line marks the ten-component threshold above which a tissue forms its own
    cross-validation stratum. (b)~Distribution of observed CLARIFi mSF2, the
    literature-median clonogenic surviving fraction at 2~Gy, with one point per independent
    component so that related cell lines count once; tissues with fewer than ten components
    are pooled. Boxes show the median and interquartile range; whiskers extend to 1.5 times
    the interquartile range. Open diamonds mark MCF-7 and DU145, two of the three lines of
    the motivating \textit{Moringa} study, shown as context only: no statement about
    \textit{Moringa} extracts or radioprotection is made or implied.}
  \label{fig:cohort}
\end{figure}

\section{Model comparison}\label{sec:comparison}

Table~\ref{tab:performance} and Figure~\ref{fig:performance} summarise the out-of-fold
performance of the five prespecified models. Two results stand out.

First, \textbf{expression features reduce the error, and the reduction is small}. Relative
to the tissue-only baseline, M2 lowers the mean absolute error by $\Delta\mathrm{MAE} =
0.00383$ mSF2 units (Equation~\ref{eq:dmae}), or 2.85\% of the baseline MAE, and it has the
lower error in \textbf{nine of ten} cross-validation repeats; the exception, repeat~7,
favours M1 by 0.0024 (Table~\ref{tab:s1}). The proliferation-adjusted sensitivity
comparison agrees in direction and size: M2P improves on M1P by 0.00387 and is better in
all ten repeats.

Second, \textbf{no model predicts well}. Out-of-sample $R^2$ ranges from $-0.018$ to
$+0.013$, so every model is approximately as accurate as predicting the cohort mean, and
the tissue-only model is marginally \emph{worse} than the global mean on both MAE and
$R^2$.

\begin{table}[tbp]
  \caption[Model comparison]{Out-of-fold performance of the five prespecified models.
    MAE, RMSE and changes in MAE are in mSF2 units; all values are component-weighted and
    averaged over ten repeats of grouped five-fold cross-validation. \enquote{Change vs M1}
    and \enquote{Change vs M0} are the reductions in MAE relative to those models (positive
    means a lower error); the paired $\Delta$MAE is the primary estimand of
    Equation~\ref{eq:dmae}, computed against M1 for M2 and against M1P for M2P.
    Descriptive only: superiority was not formally assessed. Source:
    \path{results/tables/primary_v1/publication_v1/table_2_model_comparison.csv}.}
  \label{tab:performance}
  \footnotesize\setlength{\tabcolsep}{3.2pt}
  \begin{tabular}{@{}l
      S[table-format=1.6]S[table-format=1.6]S[table-format=-1.6]
      S[table-format=-1.2]S[table-format=1.2]
      S[table-format=-1.6]S[table-format=-1.2]S[table-format=-1.6]
      S[table-format=1.6]c@{}}
    \toprule
    \thead{Model} & {\thead{MAE}} & {\thead{RMSE}} & {\thead{$R^2_{\mathrm{oos}}$}}
      & {\thead{Slope}} & {\thead{Int.}}
      & {\thead{Change}} & {\thead{\%}} & {\thead{Change}}
      & {\thead{Paired}} & \thead{Repeats} \\
    & & & & \multicolumn{2}{c}{\thead{calibration}} & \multicolumn{2}{c}{\thead{vs M1}}
      & {\thead{vs M0}} & {\thead{$\Delta$MAE}} & \thead{favouring} \\
    \midrule
    M0  & 0.133624 & 0.170755 & -0.009497 & -4.00 & 2.92 & 0.000605 & 0.45 & 0.000000 & {\na} & \na \\
    M1  & 0.134228 & 0.171494 & -0.018275 & -1.21 & 1.29 & 0.000000 & 0.00 & -0.000605 & {\na} & \na \\
    M1P & 0.134259 & 0.171304 & -0.016026 & -0.05 & 0.61 & -0.000031 & -0.02 & -0.000636 & {\na} & \na \\
    \textbf{M2} & \bfseries 0.130397 & \bfseries 0.169115 & \bfseries 0.009347 & 0.65 & 0.21
        & 0.003832 & 2.85 & 0.003227 & \bfseries 0.003832 & 9/10 \\
    M2P & 0.130388 & 0.168841 & 0.012822 & 0.63 & 0.22 & 0.003840 & 2.86 & 0.003235 & 0.003871 & 10/10 \\
    \bottomrule
  \end{tabular}
\end{table}

\begin{figure}[tbp]
  \centering
  \includegraphics[width=\linewidth,height=\ReportFigureMaxHeight,keepaspectratio]{fig_03_model_performance}
  \caption[Out-of-fold predictive performance of the five models]{Out-of-fold predictive
    performance of the five prespecified models (M0 global mean, M1 tissue only, M1P tissue
    and proliferation, M2 tissue and 50 Hallmark expression scores, M2P tissue,
    proliferation and expression). (a)~Component-weighted mean absolute error in mSF2
    units, with bars starting at zero. (b)~Out-of-sample $R^2$, where zero corresponds to
    predicting the cohort mean; every model lies close to zero. (c)~Paired difference in
    MAE in each of the ten cross-validation repeats for the primary comparison (M2 against
    M1) and the proliferation-adjusted comparison (M2P against M1P); positive values favour
    the richer model, the thick bar is the mean over repeats, and the upper axis expresses
    the difference as a percentage of MAE(M1). Repeats re-split the same 237 components, so
    their spread shows split-to-split variability, not sampling uncertainty; no interval is
    shown because no interval method was validated for this design.}
  \label{fig:performance}
\end{figure}

\paragraph{Reconciliation of the reported difference.}
Three saved files report the primary difference, and they agree exactly once the estimand
is named (Table~\ref{tab:reconciliation}). The first three entries are the same quantity
computed in three ways. The record-weighted value differs by $1.5\times10^{-5}$, 0.4\% of
the estimate; it is not a discrepancy but a different weighting, which gives components with
several records more influence. The protocol specifies the component-weighted estimate, and
that value is used throughout. The relative figure of 2.85\% is $0.0038317/0.1342282$.

\begin{table}[tbp]
  \caption{The primary difference as recorded in three saved output files.}
  \label{tab:reconciliation}
  \footnotesize
  \begin{tabular}{@{}L{0.37\linewidth}S[table-format=1.16]L{0.37\linewidth}@{}}
    \toprule
    \thead{Source} & {\thead{Value}} & \thead{Estimand} \\
    \midrule
    \path{comparisons.csv}, \path{delta_mae} & 0.0038316558867076
      & Paired, component-weighted, mean over repeats (primary) \\
    \path{model_metrics.csv}, MAE(M1)$\,-\,$MAE(M2) & 0.0038316558867076
      & Difference of the two component-weighted MAEs \\
    \path{per_repeat_delta_mae.csv}, mean of 10 repeats & 0.0038316558867076
      & Mean of the ten per-repeat values \\
    \path{comparisons.csv}, \path{record_weighted_delta_mae} & 0.0038170756763974
      & Paired, record-weighted (each of 241 records weighted equally) \\
    \bottomrule
  \end{tabular}
\end{table}

\paragraph{What the comparison establishes.}
The improvement is a descriptive estimate. Superiority was \emph{not} formally assessed: the
coverage simulation that would have validated an interval method was deferred
(Section~\ref{sec:deviation}), so no interval or $p$-value is reported. No negative-control
draws were run, so the gain cannot be placed against a permuted-outcome or random-gene-set
reference. Repeated cross-validation is not external validation: the ten repeats re-split
the same 237 components, so their spread measures split-to-split variability rather than
sampling uncertainty, and says nothing about behaviour on an independent cohort.

\section{Why out-of-sample $R^2$ stays near zero}\label{sec:shrinkage}

Figure~\ref{fig:prediction} shows the mechanism directly. Predictions from M1 and M2 occupy a
narrow band close to the cohort mean, while the observations span almost the entire
interval. The models are therefore less miscalibrated than uninformative: they shrink
heavily towards the mean because the features support little more. The calibration slopes
agree with this reading, moving from $-4.00$ for M0 and $-1.21$ for M1 to 0.65 for M2 and
0.63 for M2P (Table~\ref{tab:performance}); the strongly negative slopes of M0 and M1 are an
artefact of their predictions being nearly constant within a tissue
(Section~\ref{sec:correction}).

\begin{figure}[tbp]
  \centering
  \includegraphics[width=\linewidth,height=\ReportFigureMaxHeight,keepaspectratio]{fig_04_prediction_behaviour}
  \caption[Why the out-of-sample $R^2$ stays near zero]{Why the out-of-sample $R^2$ stays
    near zero. (a,~b)~Observed against out-of-fold predicted mSF2 for the tissue-only model
    M1 and the tissue-plus-expression model M2. Each point is one independent biological
    component (records of a component averaged; each prediction already averaged over the
    ten repeats). The dashed line is the identity, and both axes span the full range of the
    endpoint. (c)~Empirical cumulative distributions of the observed values and of the
    predictions of both models. The predictions occupy a narrow band near the cohort mean
    while the observations span almost the whole interval: the behaviour of a model with
    little usable signal rather than of a miscalibrated one. No trend line is fitted.}
  \label{fig:prediction}
\end{figure}

\section{Tissue-level and repeat-level behaviour}\label{sec:tissue}

Across the nine well-supported tissues, M2 achieves the lower error in \textbf{eight}
(Table~\ref{tab:tissue}; Figure~\ref{fig:tissue}). The largest improvements are in
Ovary/Fallopian Tube (0.0112), Breast (0.0084) and Head and Neck (0.0075). The exception is
instructive: in \textbf{Pancreas, M2 is worse by 0.0110}, nearly three times the size of the
pooled gain. A pooled improvement of this size is therefore not uniform across tissues, and
it would be misleading to present it as a general property of the model. Estimates for
tissues with fewer than ten components are shown in Table~\ref{tab:tissue} for completeness
but are unstable and are not interpreted.

Prediction stability across repeats (Figure~\ref{fig:s-variability}) is higher for the
tissue-only models, whose prediction is a shrunken tissue mean, than for the expression
models, as expected.

\begin{table}[tbp]
  \caption[Tissue-specific performance]{Component-weighted out-of-fold MAE by tissue
    (mSF2 units) and the difference $\Delta\mathrm{MAE}=\mathrm{MAE(M1)}-\mathrm{MAE(M2)}$;
    positive values mean the expression model had the lower error. Tissues below the rule
    have fewer than ten components; their estimates are descriptive only and unstable.
    Source: \path{results/tables/primary_v1/publication_v1/table_3_tissue_performance.csv}.}
  \label{tab:tissue}
  \small
  \begin{tabular}{@{}l S[table-format=2.0] S[table-format=1.4] S[table-format=1.4]
      S[table-format=1.4] S[table-format=-1.4] l@{}}
    \toprule
    \thead{Tissue} & {\thead{Components}} & {\thead{MAE M0}} & {\thead{MAE M1}}
      & {\thead{MAE M2}} & {\thead{$\Delta$MAE}} & \thead{Direction} \\
    \midrule
    Lung & 46 & 0.1251 & 0.1258 & 0.1249 & 0.0009 & M2 lower \\
    CNS/Brain & 23 & 0.1558 & 0.1580 & 0.1538 & 0.0042 & M2 lower \\
    Esophagus/Stomach & 23 & 0.1252 & 0.1279 & 0.1277 & 0.0002 & M2 lower \\
    Bowel & 21 & 0.1512 & 0.1511 & 0.1471 & 0.0040 & M2 lower \\
    Head and Neck & 18 & 0.0756 & 0.0720 & 0.0645 & 0.0075 & M2 lower \\
    Breast & 17 & 0.1051 & 0.1050 & 0.0966 & 0.0084 & M2 lower \\
    Pancreas & 16 & 0.0921 & 0.0934 & 0.1044 & -0.0110 & \textbf{M2 higher} \\
    Bladder/Urinary Tract & 11 & 0.1252 & 0.1255 & 0.1202 & 0.0053 & M2 lower \\
    Ovary/Fallopian Tube & 10 & 0.2216 & 0.2272 & 0.2160 & 0.0112 & M2 lower \\
    \midrule
    Skin & 9 & 0.1601 & 0.1630 & 0.1653 & -0.0023 & M2 higher \\
    Bone & 8 & 0.1287 & 0.1317 & 0.1290 & 0.0027 & M2 lower \\
    Uterus & 8 & 0.1866 & 0.1791 & 0.1484 & 0.0308 & M2 lower \\
    Pleura & 6 & 0.1855 & 0.1802 & 0.1788 & 0.0014 & M2 lower \\
    Prostate & 5 & 0.1078 & 0.1102 & 0.0828 & 0.0275 & M2 lower \\
    Kidney & 3 & 0.1627 & 0.1674 & 0.1666 & 0.0008 & M2 lower \\
    Liver & 3 & 0.0430 & 0.0435 & 0.0696 & -0.0261 & M2 higher \\
    Peripheral Nervous System & 3 & 0.2041 & 0.2036 & 0.1985 & 0.0051 & M2 lower \\
    Soft Tissue & 3 & 0.1117 & 0.1109 & 0.1255 & -0.0146 & M2 higher \\
    Thyroid & 3 & 0.1556 & 0.1576 & 0.1306 & 0.0270 & M2 lower \\
    Biliary Tract & 1 & 0.4168 & 0.4146 & 0.3728 & 0.0418 & M2 lower \\
    \bottomrule
  \end{tabular}
\end{table}

\begin{figure}[tbp]
  \centering
  \includegraphics[width=\linewidth,height=\ReportFigureMaxHeight,keepaspectratio]{fig_05_tissue_performance}
  \caption[Performance by tissue of origin]{Performance by tissue of origin for the nine
    tissues with at least ten independent components; the number of components is given
    beside each tissue name. (a)~Component-weighted out-of-fold MAE of each model, in mSF2
    units; lighter cells are lower error. (b)~The difference MAE(M1)$\,-\,$MAE(M2) for the
    same tissues in the same order; positive values (blue) mean the expression model had
    the lower error. Differences are small relative to the errors themselves and
    descriptive only: they come from one cross-validated run and no per-tissue inference
    was performed.}
  \label{fig:tissue}
\end{figure}

\section{A reporting correction}\label{sec:correction}

The first version of the within-tissue Spearman analysis reported strongly negative rank
correlations for M1 in all nine larger tissues. Inspection of the code established the
cause. Every fitted M1 predicts a single value per tissue, the shrunken training-fold tissue
mean, so a held-out line's pooled out-of-fold prediction is effectively the mean of the
\emph{other} lines of its tissue, which falls as the line's own value rises. This is a
leave-fold-out mean artefact, not a biological finding, and for M0 and M1 the statistic
measures no within-tissue ranking at all. The corrected table (Table~\ref{tab:s4}) reports
the statistic as undefined wherever a model's prediction is constant within a tissue, and
its defined values carry an explicit downward-bias caveat. The original pooled within-tissue
correlations are not interpreted anywhere in this report; the first version is retained
unchanged for audit.

\chapter{Robustness and tissue-generalisation analyses}\label{ch:robustness}

The results of Chapter~\ref{ch:results} come from the frozen primary benchmark. A separate
and clearly labelled layer of analyses tests how durable its small advantage is, without
altering it (Table~\ref{tab:robustness}). Three of the four planned analyses were executed;
the negative controls were configured but not run.

\begin{table}[tbp]
  \caption{The exploratory and robustness layer: questions, runs and execution status.}
  \label{tab:robustness}
  \footnotesize
  \begin{tabular}{@{}L{0.20\linewidth}L{0.41\linewidth}L{0.13\linewidth}L{0.19\linewidth}@{}}
    \toprule
    \thead{Analysis} & \thead{Question} & \thead{Run} & \thead{Status} \\
    \midrule
    Random forest M3 & Does a different model family extract more signal from the same
      features and splits? & \path{exploratory_m3_v1} & Completed, 200/200 tasks \\
    Leave-one-tissue-out & Does an expression-only model transfer to a tissue it has never
      seen? & \path{loto_v1} & Completed, 18/18 tasks \\
    Outcome-scale sensitivity & Does the advantage survive the neglog and logit scales, or
      is it a raw-scale artefact? & \path{scales_v1} & Completed, 400/400 tasks \\
    Negative controls (1,000 draws of each null) & Is the advantage larger than the pipeline
      produces when expression cannot carry signal? Are any gene sets more important than
      matched random sets? & \path{controls_v1}
      & Configured, about 10--20~h at six workers; \textbf{not executed} \\
    \bottomrule
  \end{tabular}
\end{table}

\section{A different model family does not do better}\label{sec:m3}

The random forest M3, fitted on the same splits with the same nested tuning, reached MAE
0.13148 and out-of-sample $R^2$ $+0.00643$, against 0.13040 and $+0.00935$ for M2
(Table~\ref{tab:r1}). Its gain over the tissue-only baseline, $\Delta\mathrm{MAE}=0.00275$,
is \textbf{smaller} than the 0.00383 of M2. \textbf{M3 therefore performs worse than M2, and
that is the result reported.} Within the same run, M0, M1 and M2 reproduced the primary run
to the last decimal place, which also confirms that the frozen benchmark is reproducible.

\begin{table}[tbp]
  \caption[Exploratory model comparison including the random forest M3]{Exploratory model
    comparison including the random forest M3, on the raw mSF2 scale (MAE and RMSE in mSF2
    units). M0, M1 and M2 reproduce Table~\ref{tab:performance}. Source:
    \path{results/tables/exploratory_m3_v1/model_metrics.csv}; the paired differences
    quoted in the text are from \path{comparisons.csv} in the same folder.}
  \label{tab:r1}
  \small
  \begin{tabular}{@{}l S[table-format=1.6] S[table-format=1.6] S[table-format=-1.6]
      S[table-format=-1.2] S[table-format=1.2]@{}}
    \toprule
    \thead{Model} & {\thead{MAE}} & {\thead{RMSE}} & {\thead{$R^2_{\mathrm{oos}}$}}
      & {\thead{Calibration slope}} & {\thead{Calibration intercept}} \\
    \midrule
    M0 & 0.133624 & 0.170755 & -0.009497 & -4.00 & 2.92 \\
    M1 & 0.134228 & 0.171494 & -0.018275 & -1.21 & 1.29 \\
    M2 & 0.130397 & 0.169115 & 0.009347 & 0.65 & 0.21 \\
    M3 & 0.131480 & 0.169394 & 0.006434 & 0.55 & 0.26 \\
    \bottomrule
  \end{tabular}
\end{table}

\section{Transfer to an unseen tissue is weak and conditional}\label{sec:loto}

Holding out each of the nine well-supported tissues in turn (185 components, 187 records),
the expression-only model M2X reached MAE 0.12267 and out-of-sample $R^2$ \textbf{+0.0308},
against 0.12753 and $-0.0084$ for the global mean ($\Delta\mathrm{MAE}=0.00487$ mSF2 units).
This is the only analysis in the project in which out-of-sample $R^2$ is clearly above zero.
It is, however, a different and easier comparison than the primary one --- expression alone
against a global mean, with tissue removed from both sides --- so it is not comparable with
the primary 0.00383 and does not revise it.

The pooled figure conceals an uneven picture (Table~\ref{tab:r2}). The gain is largest when
Ovary/Fallopian Tube ($+0.0156$) or Head and Neck ($+0.0153$) is held out, followed by
Breast ($+0.0119$), Bowel ($+0.0092$), Bladder/Urinary Tract ($+0.0079$), Lung ($+0.0049$)
and CNS/Brain ($+0.0025$). It is \textbf{negative for two tissues, Pancreas ($-0.0142$) and
Esophagus/Stomach ($-0.0030$)}, where the expression-only model predicted worse than the
global mean. The Pancreas loss is nearly three times the pooled gain, and Pancreas is also
the one well-supported tissue in which M2 was worse than M1 in the primary benchmark
(Table~\ref{tab:tissue}). Transfer of expression information to an unseen tissue is
therefore conditional on the tissue, and the tissue-level results, not only the pooled
gain, belong in any statement about generalisation.

\begin{table}[tbp]
  \caption[Leave-one-tissue-out performance by held-out tissue]{Leave-one-tissue-out
    performance of the expression-only model M2X by held-out tissue. $\Delta$MAE is
    MAE(global mean)$\,-\,$MAE(M2X) in mSF2 units, so positive values favour expression;
    the Spearman correlation is that of M2X within the held-out tissue. The pooled row is
    from \path{comparisons.csv} and \path{model_metrics.csv} of the same run. Source:
    \path{results/tables/loto_v1/held_out_tissue.csv}.}
  \label{tab:r2}
  \small
  \begin{tabular}{@{}l S[table-format=3.0] S[table-format=-1.4] S[table-format=-1.4]
      S[table-format=-1.3]@{}}
    \toprule
    \thead{Held-out tissue} & {\thead{Components}} & {\thead{$\Delta$MAE}}
      & {\thead{Mean prediction shift}} & {\thead{Spearman $\rho$}} \\
    \midrule
    Bladder/Urinary Tract & 11 & 0.0079 & -0.0425 & 0.327 \\
    Bowel & 21 & 0.0092 & 0.0096 & 0.231 \\
    Breast & 17 & 0.0119 & 0.0290 & 0.324 \\
    CNS/Brain & 23 & 0.0025 & 0.0735 & 0.359 \\
    Esophagus/Stomach & 23 & \bfseries -0.0030 & 0.0273 & -0.124 \\
    Head and Neck & 18 & 0.0153 & -0.0284 & 0.187 \\
    Lung & 46 & 0.0049 & -0.0221 & 0.178 \\
    Ovary/Fallopian Tube & 10 & 0.0156 & -0.0452 & 0.576 \\
    Pancreas & 16 & \bfseries -0.0142 & 0.0115 & -0.062 \\
    \midrule
    Pooled (nine tissues) & 185 & 0.0049 & {\na} & {\na} \\
    \bottomrule
  \end{tabular}
\end{table}

\section{The advantage depends on the outcome scale}\label{sec:scales}

The most important negative finding of the robustness layer concerns the outcome scale.
Repeating the M2-versus-M1 comparison with 20 repeats on two transformed scales gave
$\Delta\mathrm{MAE}=0.00023$ on the neglog scale --- essentially nothing, with the
out-of-sample $R^2$ of M2 falling to $-0.0322$ --- and 0.00980 on the logit scale
(Table~\ref{tab:r3}). The direction favours M2 on all three scales. The magnitudes, however,
are expressed in the units of each transformed scale, not in mSF2, so they cannot be
compared directly with the raw-scale 0.00383; each must be read against the baseline error
on its own scale, which is also given in the table. On the neglog scale the advantage is
negligible; on the logit scale the out-of-sample $R^2$ of M2 is slightly positive (0.0045)
while that of M1 is negative. The raw scale
remains the protocol's prespecified primary scale, and these results are reported beside it
rather than replacing it. Together they mean that the primary estimate should be read as
\textbf{small, consistent in direction and sensitive to how the outcome is expressed}.

\begin{table}[tbp]
  \caption[Outcome-scale sensitivity of the primary comparison]{Outcome-scale sensitivity of
    the primary comparison (M2 against M1). $\Delta$MAE and MAE are in the units of each
    scale and are not comparable across scales. The raw row is the primary run (10 repeats);
    the transformed scales used 20 repeats on the same 237 components and 241 records.
    Descriptive only; no confirmatory method was adopted. Sources:
    \path{results/tables/scales_v1/comparisons.csv} and \path{model_metrics.csv};
    raw scale from Table~\ref{tab:performance}.}
  \label{tab:r3}
  \small
  \begin{tabular}{@{}l S[table-format=1.5] S[table-format=1.5] S[table-format=2.0]
      S[table-format=1.4] S[table-format=-1.4] S[table-format=-1.4]@{}}
    \toprule
    \thead{Scale} & {\thead{$\Delta$MAE}} & {\thead{Record-weighted}} & {\thead{Repeats}}
      & {\thead{MAE M1}} & {\thead{$R^2_{\mathrm{oos}}$ M1}} & {\thead{$R^2_{\mathrm{oos}}$ M2}} \\
    \midrule
    raw (primary) & 0.00383 & 0.00382 & 10 & 0.1342 & -0.0183 & 0.0093 \\
    neglog & 0.00023 & 0.00017 & 20 & 0.2596 & -0.0164 & -0.0322 \\
    logit & 0.00980 & 0.00957 & 20 & 0.6398 & -0.0142 & 0.0045 \\
    \bottomrule
  \end{tabular}
\end{table}

\section{Rules governing the exploratory layer}\label{sec:rules}

Four rules keep this layer separate from the primary benchmark.

\paragraph{It cannot promote a new primary model.} M3 was configured \emph{after} the
primary results were seen, so M2 remains the primary benchmark whatever M3 scores. In the
event M3 scored worse, which is the expected outcome for a tree ensemble on 237 components
and 50 coarse features and suggests that the ceiling is set by the data rather than by the
model family.

\paragraph{Consistency with the frozen run is verified mechanically.} A dedicated check
confirms that an exploratory run covers the same 241 records and 237 components; that no
biological component straddles a training and test fold or two inner folds; that the split
table and its checksum are identical to those of the primary run; that the input
checksums, base seed and fold counts match; that tuning is nested; that models shared with
the primary run reproduce its predictions exactly; and that the primary run's own files
still match their recorded checksums.

\paragraph{Negative controls would not by themselves create inference.} With 1,000 draws,
the gene-set $p$-values that are currently empty would become available, and the observed
$\Delta\mathrm{MAE}$ could be placed descriptively against a permuted-outcome null. Tiers A
and B would stay closed, because that gate also requires a formally demonstrated prediction
gain, which needs the interval method that the deferred coverage simulation was designed to
validate. A permutation test would be a legitimate alternative route, but adopting it
requires a protocol amendment recorded \emph{before} the null distribution is inspected;
written afterwards, it would be a post-hoc test chosen in sight of the result and would have
to be reported as such.

\paragraph{The within-tissue correlation artefact is avoided.} None of these analyses
reports the pooled within-tissue rank correlation produced by the first version of that
table; the corrected statistic is undefined wherever a model predicts a constant within a
tissue (Section~\ref{sec:correction}).

\chapter{Explainability and pathway prioritisation}\label{ch:explainability}

\section{Gene-set importance}\label{sec:importance}

All 50 fits of M2 completed and were analysed. Ranked by mean held-out permutation importance
(Equation~\ref{eq:importance}; Figure~\ref{fig:importance}a; Table~\ref{tab:s2}), the leading
gene set is \textbf{HALLMARK\_P53\_PATHWAY}, with mean importance 0.00343 mSF2 units,
importance above zero in 90\% of fits, sign consistency 0.92 and a predominantly negative
coefficient. It is followed by COMPLEMENT (0.00118), OXIDATIVE\_PHOSPHORYLATION (0.00107),
UV\_RESPONSE\_UP (0.00100), MYC\_TARGETS\_V2, KRAS\_SIGNALING\_UP and UV\_RESPONSE\_DN.

Three observations matter more than the ranking itself.
\begin{itemize}
  \item \textbf{The magnitudes are tiny.} The largest mean importance, 0.0034 mSF2 units, is
    of the same order as the entire predictive gain, so the ranking describes a very small
    effect.
  \item \textbf{Stability is modest.} Across the 50 fits, the mean Kendall $\tau$ between
    importance rankings is 0.225, the mean rank-biased overlap 0.518 and the mean top-10
    Jaccard index 0.391: the same gene sets do not reliably appear at the top of every fit.
  \item \textbf{Correlated sets share signal.} Several ranked sets are correlated with
    another set at $\lvert\rho\rvert\ge0.80$ (Figure~\ref{fig:s-correlation}), including
    COMPLEMENT with APOPTOSIS and MYC\_TARGETS\_V2 with MYC\_TARGETS\_V1, so importance may
    be shared or exchanged between them rather than attributable to one.
\end{itemize}

Equally important is what the ranking does \emph{not} contain. \textbf{DNA repair, G2M
checkpoint and hypoxia, the pathways a radiobiologist would expect, were not prioritised}:
G2M checkpoint ranked 13th with importance above zero in only 8\% of fits, hypoxia 10th with
10\%, and DNA repair 47th. Reporting this guards against presenting a familiar
radiation-response narrative that the data do not support.

\begin{figure}[tbp]
  \centering
  \includegraphics[width=\linewidth,height=\ReportFigureMaxHeight,keepaspectratio]{fig_06_gene_set_importance}
  \caption[Exploratory importance of the Hallmark gene-set scores in M2]{Exploratory
    importance of the Hallmark gene-set scores in the expression model M2. (a)~The fifteen
    highest-ranked gene sets. Grey points are the grouped permutation importance in each of
    the fifty model fits (ten repeats by five folds), computed on held-out rows only as the
    increase in component-weighted MAE when a score is permuted across components; diamonds
    are the means over fits, coloured by the prespecified tier. The spread across fits is
    split-to-split variability, not a confidence interval. (b)~The two tier-C criteria for
    all fifty gene sets: the share of fits with importance above zero and the share of fits
    in which the set received a non-zero coefficient; the shaded corner satisfies both
    thresholds (dashed lines at 0.50). Tiers A and B were unreachable by design. Shares are
    proportions of overlapping cross-validation fits, not probabilities. The rankings
    describe one fitted model family on one dataset; they are not validated biomarkers,
    causal mechanisms, radiation-response pathways or molecular effect sizes.}
  \label{fig:importance}
\end{figure}

\section{The pathway shortlist}\label{sec:pathway-shortlist}

Seven of the 50 gene sets reached \textbf{tier~C} and 43 were not prioritised; \textbf{no set
reached tier A or B}, because both require a formally demonstrated prediction gain and at
least 1,000 negative-control draws, and the run has neither (Figure~\ref{fig:importance}b).
The gene-set tests record zero completed draws of both the permuted-outcome and the
random-gene-set nulls, so attribution is not permitted and the stability gate cannot be
evaluated.

The shortlist has a precise meaning. Its members are the gene-set scores whose permutation
importance was positive in at least half of the cross-validation fits and which received a
non-zero coefficient in at least half, within one model family on one dataset. It is not a
gene-level finding, a validated biomarker, a causal mechanism, a radiation-response pathway
or a molecular effect size, and it is not evidence that modulating any of these pathways
would change radiosensitivity.

\section{Cell lines the model reproduces least well}\label{sec:discrepancies}

Of the 241 measured lines, 12 reached discrepancy tier~A, nine tier~B and 27 tier~C; 193 were
not prioritised (Figure~\ref{fig:discrepancies}; Table~\ref{tab:s3}). The highest-ranked
discrepancies are dominated by lines with extreme observed values --- HT-144 (observed 0.08,
predicted 0.59 under M2), KNS-42 (0.125 against 0.60), HEC-1-B (0.97 against 0.56) and
D283~Med (0.18 against 0.58) --- which is the shrinkage of Figure~\ref{fig:prediction} seen one line at a time
rather than a set of biologically unusual cell lines.

The triage separates possible explanations. Several tier A and B lines carry technical flags
(small tissue group, related lines in the cohort, a Cellosaurus problem flag on a relative)
and are routed to technical review first; the remainder are labelled as unexplained
discrepancies requiring follow-up. Because the tier thresholds are cohort percentiles, a
roughly fixed fraction of lines is always flagged: a flag means \enquote{relatively large in
this cohort}, not \enquote{abnormal}.

Of the two context lines from the motivating study, \textbf{MCF-7 was not prioritised}: its
held-out errors lie within the bulk of the cohort (observed 0.52 against 0.574 under M1 and
0.534 under M2). \textbf{DU145 reached tier~C} with four technical flags --- stranded RNA
profile, small tissue group, unstable prediction across repeats and strong model
disagreement --- and was observed at 0.71 against predictions of 0.580 under M1 and 0.650
under M2. Both tiers were assigned by the same rule as for every other line. No individual
prediction error is evidence of radioprotection, radiosensitisation or any causal mechanism,
and lines selected for extreme residuals are expected to agree better on re-measurement even
if nothing is wrong, through regression to the mean.

\begin{figure}[tbp]
  \centering
  \includegraphics[width=\linewidth,height=\ReportFigureMaxHeight,keepaspectratio]{fig_07_cell_line_discrepancies}
  \caption[Cell lines whose radiosensitivity the model reproduces least well]{Cell lines whose
    measured radiosensitivity the model reproduces least well. (a)~Residual (observed minus
    predicted mSF2 under M2, predictions averaged over repeats) against the observed value
    for every measured line, coloured by prespecified discrepancy tier; triangles mark lines
    with at least one technical flag. The strong dependence of the residual on the observed
    value is shrinkage: because the models predict close to the cohort mean, lines with
    extreme measured values necessarily have the largest residuals. Circled points are MCF-7
    and DU145, labelled for their relevance to the motivating study; their tiers are not
    altered. (b)~Flags recorded for the tier A and B lines, in rank order, separating flags
    about prediction behaviour (instability across repeats, disagreement between M1 and M2)
    from flags about the data or cell-line identity (stranded RNA profile, related lines in
    the cohort, a Cellosaurus problem flag on a relative, small tissue group). These are
    measured lines used in model development, not independent validation samples; a large
    residual is a signal to re-check identity, provenance and the original measurement, not
    evidence of unusual radiobiology.}
  \label{fig:discrepancies}
\end{figure}

\chapter{\textit{Moringa} compound evidence and candidate prioritisation}\label{ch:compounds}

\section{Verified constituents of \textit{Moringa oleifera} seed and leaf}\label{sec:constituents}

Curation produced a provenance-carrying inventory of \textit{M.~oleifera} seed and leaf
constituents drawn from primary analytical studies, with the extraction solvent,
identification method, DOI verification route and reading depth of each report recorded
(Table~\ref{tab:c1}). Three findings are substantive in themselves.

\paragraph{Only one constituent is confirmed in the solvent the original study used.}
\textcite{abdelsayed2021jfda} isolated \textbf{niazimicin}, a thiocarbamate glycoside, from
the \textbf{seed methylene-chloride fraction}, confirmed its structure by HRMS, UV, IR, EI-MS
and one- and two-dimensional NMR, and quantified it by validated UPLC-MS/MS in both the seed
and the leaf \ce{CH2Cl2} fractions. The same study did \emph{not} detect niazirin or
glucomoringin there. The only other dichloromethane-specific identifications are GC-MS
spectral-library matches in leaf DCM material, read at abstract level:
4-hydroxyphenylacetonitrile and methyl 4-hydroxyphenylacetate \parencite{mohdfisall2021},
and palmitic acid and \textit{cis}-vaccenic acid \parencite{khandelwal2019}. The latter
appeared in a journal indexed in neither PubMed nor Europe PMC, its full compound list could
not be accessed, and electron-impact library matching cannot distinguish the $\Delta^{9}$
and $\Delta^{11}$ C18:1 isomers without authentic standards.

\paragraph{The identification of several widely cited \enquote{Moringa compounds} is weaker
than their reputation.} Free quercetin and kaempferol rest only on tentative LC-HRMS
dereplication without authentic standards and may form by hydrolysis of the abundant
glycosides during extraction; the forms that primary analytical work actually establishes in
leaf are the \textbf{glycosides} isoquercitrin, astragalin and rutin. Under the prespecified
rules the tentative aglycone reports still count as primary seed or leaf reports, so
quercetin and kaempferol \emph{are} assessable, but only in the weakest identification band
($D_1=1$). The pharmacokinetic caution applies in both directions: aglycone evidence may not
be transferred to a glycoside, or the reverse, without explicit justification. Benzyl
isothiocyanate is \textbf{not} a documented seed or leaf constituent: its precursor, benzyl
glucosinolate, is reported in \emph{roots}, while seed and leaf carry the
rhamnosyloxy-benzyl glucosinolate instead. Gallic acid and caffeic acid yielded no primary
report in \textit{M.~oleifera} seed or leaf and do not appear in the inventory.

\paragraph{Identifier hygiene was not a formality.} Four public records returned by a name
search are wrong or ambiguous for this purpose, and using them would have propagated an
error: a mis-structured carbamate record (\ce{C16H23NO8}, with the ethoxy group on the
sugar); a \enquote{niaziminin A} record whose structure is not the thiocarbamate; a
niazimicin record without stereochemistry; and glucomoringin, for which every available
record is a potassium salt with a different stereo-InChIKey and none could be matched to the
natural structure. The glucomoringin identifier remains \textbf{unresolved}.

\section{The candidate shortlist}\label{sec:shortlist}

The evidence linking these constituents to radiation-response biology was curated
separately, from primary studies of purified compounds only (Table~\ref{tab:c2}), and the
prespecified score of Equation~\ref{eq:priority} was then applied
(Table~\ref{tab:ranking}; Figures~\ref{fig:priority} and~\ref{fig:evidence-map}).
\textbf{Eighteen compounds entered the ranked shortlist and five reached tier~A}; three
further constituents scored below the tier-C threshold and were not ranked, and benzyl
isothiocyanate was not assessable. The written justification of every point is given in
Table~\ref{tab:c3reasons}.

\begin{table}[tbp]
  \caption[Candidate ranking]{Candidate ranking of documented \textit{M.~oleifera} seed and
    leaf constituents under the prespecified four-dimension score (Equation~\ref{eq:priority};
    $D_1$ identity 0--2, $D_2$ radiation biology 0--3, $D_3$ independent sources 0--2, $D_4$
    feasibility 0--1). Extraction evidence is shown beside the score and is not part of it.
    Rows follow the ranking file. Every tier is exploratory and hypothesis-generating.
    Source: \path{results/tables/primary_v1/compounds_v1/table_c3_candidate_ranking.csv}.}
  \label{tab:ranking}
  \footnotesize\setlength{\tabcolsep}{3.4pt}
  \begin{tabular}{@{}r L{0.235\linewidth} l l c c c c c L{0.10\linewidth} L{0.15\linewidth}@{}}
    \toprule
    \thead{Rank} & \thead{Compound} & \thead{Part} & \thead{Extraction}
      & \thead{$D_1$} & \thead{$D_2$} & \thead{$D_3$} & \thead{$D_4$} & \thead{Total}
      & \thead{Tier} & \thead{Liability} \\
    \midrule
    1 & $\alpha$-Tocopherol & both & other solvent & 2 & 3 & 2 & 1 & 8 & \textbf{A} & \na \\
    2 & Chlorogenic acid & leaf & other solvent & 1 & 3 & 2 & 1 & 7 & \textbf{C} (capped from A)
      & protects tumour cells (CGA-01) \\
    3 & Kaempferol & both & indirect & 1 & 3 & 2 & 1 & 7 & \textbf{A} & \na \\
    4 & Quercetin & both & indirect & 1 & 3 & 2 & 1 & 7 & \textbf{A} & \na \\
    5 & Rutin & leaf & other solvent & 1 & 3 & 2 & 1 & 7 & \textbf{A} & \na \\
    6 & $\beta$-Sitosterol & seed & other solvent & 1 & 3 & 2 & 1 & 7 & \textbf{A} & \na \\
    \na & Benzyl isothiocyanate & other & indirect & 0 & 3 & 2 & 1 & 6 & not assessable & \na \\
    7 & Niazirin & both & other solvent & 2 & 1 & 2 & 1 & 6 & B & \na \\
    8 & Oleic acid & seed & other solvent & 1 & 2 & 1 & 1 & 5 & B & \na \\
    9 & Astragalin & leaf & other solvent & 1 & 1 & 2 & 1 & 5 & C (capped from B)
      & genotoxic (AST-01) \\
    10 & Moringin & both & other solvent & 1 & 1 & 2 & 1 & 5 & B & \na \\
    11 & Phytol & leaf & other solvent & 1 & 1 & 2 & 1 & 5 & B & \na \\
    12 & Glucomoringin & both & other solvent & 2 & 0 & 2 & 0 & 4 & C & \na \\
    13 & Niazimicin & both & \textbf{DCM} & 2 & 0 & 2 & 0 & 4 & C & \na \\
    14 & Acetylated moringin\textsuperscript{a} & leaf & other solvent & 2 & 0 & 1 & 0 & 3 & C & \na \\
    15 & Niaziminin & leaf & other solvent & 2 & 0 & 1 & 0 & 3 & C & \na \\
    16 & \textit{O}-Ethyl rhamnosyloxybenzyl carbamate\textsuperscript{b} & seed & other solvent
      & 2 & 0 & 1 & 0 & 3 & C & \na \\
    17 & Isoquercitrin & leaf & other solvent & 1 & 0 & 1 & 0 & 2 & C & \na \\
    18 & Palmitic acid & both & \textbf{DCM} & 1 & 0 & 1 & 0 & 2 & C & \na \\
    \na & 4-Hydroxyphenylacetonitrile & both & \textbf{DCM} & 1 & 0 & 0 & 0 & 1 & not prioritised & \na \\
    \na & Methyl 4-hydroxyphenylacetate & leaf & \textbf{DCM} & 1 & 0 & 0 & 0 & 1 & not prioritised & \na \\
    \na & \textit{cis}-Vaccenic acid & leaf & \textbf{DCM} & 1 & 0 & 0 & 0 & 1 & not prioritised & \na \\
    \bottomrule
  \end{tabular}
  \tablenote{Extraction: \textbf{DCM}, identified in a dichloromethane extract or fraction of
    seed or leaf; other solvent, identified in seed or leaf with another solvent; indirect,
    identification indirect or unconfirmed. For benzyl isothiocyanate the plant part is
    \enquote{other} because its only \textit{Moringa} report concerns its precursor in roots.
    \textsuperscript{a}\,4-[(4$'$-\textit{O}-Acetyl-$\alpha$-L-rhamnosyloxy)benzyl]
    isothiocyanate. \textsuperscript{b}\,\textit{O}-Ethyl-4-($\alpha$-L-rhamnosyloxy)benzyl
    carbamate.}
\end{table}

Three features of the ranking matter more than the order itself.

\paragraph{The liability cap did real work.} Chlorogenic acid scored 7 of 8, equal to four
tier-A compounds, and is nevertheless held at tier~C because one primary study reports that
it \emph{protects hepatocellular carcinoma cells} from radiation. A score alone would have
promoted it; the prespecified cap prevented that.

\paragraph{The exclusion rule did real work.} Benzyl isothiocyanate scored 6 on biology and
sources and is still \emph{not assessable}, because it has no primary report as a seed or leaf
constituent. It cannot be ranked however strong its radiation evidence is.

\paragraph{The mismatch exposed by the curation is visible in the ranking.} Niazimicin holds
the only full identity score among compounds confirmed in a dichloromethane fraction
($D_1=2$) and scores \emph{zero} on radiation biology, so it sits at tier~C. Moringin, the
signature \textit{Moringa} isothiocyanate, reaches only tier~B. Four of the five tier-A
compounds hold $D_1=1$, a spectral-library or dereplication match, and none of the five has
been found in a dichloromethane extract.

\section{What the published biological evidence shows}\label{sec:biology}

The curated evidence establishes the following.
\begin{itemize}
  \item \textbf{Moringin, the isothiocyanate most characteristic of \textit{Moringa}, has no
    ionising-radiation study.} Exhaustive searching found no primary study combining
    moringin or glucomoringin with ionising radiation. Its evidence is confined to work
    without radiation: pro-apoptotic effects involving p53, CDKN1A, BAX and caspases in
    neuroblastoma cells \parencite{cirmi2019}, G2/M arrest in prostate cancer cells
    \parencite{abdkarim2023}, NRF2 activation in a motor-neuron model
    \parencite{gugliandolo2024}, and NF-$\kappa$B inhibition with a myeloma xenograft
    \parencite{brunelli2010}. Because chemical-class analogy earns no points, moringin gains
    nothing from the benzyl isothiocyanate radiation literature, and that analogy is
    treated as an untested hypothesis.
  \item \textbf{The best-evidenced radiation candidates are not \textit{Moringa}-specific.}
    $\beta$-Sitosterol has normal-tissue radioprotection evidence from two independent
    groups \parencite{li2007,elbakary2026}, as does rutin \parencite{patil2017,thabet2018},
    although one research group contributes to both compounds' evidence;
    $\alpha$-tocopherol has a rat crypt-microcolony clonogenic result
    \parencite{felemovicius1995} and a randomised controlled trial in radiation mucositis
    \parencite{ferreira2004}. All three are documented in \textit{Moringa} seed or leaf, but
    \textbf{none has been found in a dichloromethane extract}, which is why the extraction
    evidence is shown beside the score rather than within it.
  \item \textbf{One candidate carries an explicit liability.} Chlorogenic acid has
    normal-tissue protection evidence in bone marrow \parencite{abraham1993} and testis
    \parencite{abedpour2022}, but \textcite{yin2022} report that it \emph{protects
    hepatocellular carcinoma cells} from radiation through NRF2 and conclude that it could
    hinder radiotherapy. Under the prespecified rule this caps its tier at C. Astragalin is
    capped for a related reason: its only DNA-damage-relevant primary signal is
    \emph{pro-genotoxic} \parencite{saito2012}, and the one astragalin \enquote{radiation}
    paper concerns UVB, which is not ionising radiation.
  \item \textbf{An informative negative was retained.} Membrane oleate enrichment did
    \textbf{not} alter the clonogenic radiosensitivity of human lymphoid cells
    \parencite{george1983}. This direct clonogenic test argues against prioritising oleic
    acid despite its abundance in seed oil.
\end{itemize}

\begin{figure}[p]
  \centering
  \includegraphics[width=\linewidth,height=0.62\textheight,keepaspectratio]{fig_08_compound_priority}
  \caption[Evidence-based prioritisation of \textit{Moringa} seed and leaf
    constituents]{Evidence-based prioritisation of documented \textit{Moringa oleifera} seed
    and leaf constituents for future experimental work. (a)~Points awarded in each of the
    four prespecified dimensions (numbers are points; shading is the fraction of the
    dimension's maximum), beside the extraction-specific evidence: identified in a
    dichloromethane extract or fraction of seed or leaf (dark), in seed or leaf with another
    solvent (light), or only indirectly (white). (b)~Total priority score out of eight,
    coloured by tier. The extraction column is kept outside the score because the original
    study used dichloromethane extracts whose composition was never measured. The scoring
    rules were fixed before any score was computed; every tier, including A, is
    hypothesis-generating, and no tier asserts radioprotection, radiosensitisation, potency
    or any effect of the original extracts. Two systematic names are shortened in the
    figure (see the note to Table~\ref{tab:ranking}).}
  \label{fig:priority}
\end{figure}

\paragraph{Reading Figure~\ref{fig:priority}.} Down the extraction column and across the
four dimensions, the \textit{Moringa}-specific compounds with the darkest identity cell ---
niazimicin, glucomoringin, niazirin, niaziminin, the acetylated isothiocyanate and the
carbamate glycoside, all with a structurally confirmed presence in seed or leaf --- are
precisely those whose radiation-biology cell is empty or nearly so. The one compound with
full scores for both identity and radiation biology is $\alpha$-tocopherol, a ubiquitous
plant metabolite that has not been found in a dichloromethane extract; the other four
tier-A compounds carry only one identity point. Panel~(b) shows the effect of the liability
rule: chlorogenic acid's total of 7 equals that of four tier-A compounds, yet its bar is
coloured as tier~C. A reader comparing bar lengths alone would place it among the leaders;
the prespecified cap prevents that, and the figure shows the discrepancy rather than hiding
it.

\begin{figure}[p]
  \centering
  \includegraphics[width=\linewidth,height=0.62\textheight,keepaspectratio]{fig_09_compound_evidence_map}
  \caption[What the published evidence links each shortlisted compound to]{What the
    published evidence links each shortlisted compound to. (a)~Radiation-response processes
    for which a primary study of the purified compound exists: symbol size is the number of
    independent research groups, and colour distinguishes studies using ionising radiation
    from studies of the same process without radiation. Compounds are in priority order. An
    empty row means that no qualifying primary study was located, which is a gap in the
    literature, not evidence of no effect. (b)~Number of genes measured in a published
    experiment on the compound that are also members of a Hallmark gene set ranked by the
    exploratory explainability analysis (Table~\ref{tab:c4}); blank cells mean no overlap was
    recorded. An overlap marks a point where a compound and a model-highlighted gene set
    touch the same genes. It is a hypothesis to be tested, not evidence that the compound
    modifies the gene set, the pathway or clonogenic SF2, and the gene-set ranking itself is
    capped at tier~C. Reported directions of effect are given in Table~\ref{tab:c2} and are
    never merged.}
  \label{fig:evidence-map}
\end{figure}

\paragraph{Reading Figure~\ref{fig:evidence-map}.} Attributed compound by compound, the
biological evidence is thin. Only a minority of cells in panel~(a) are filled, and several
shortlisted compounds have no ionising-radiation evidence on their row: moringin, the most
\textit{Moringa}-characteristic candidate, appears only in the non-radiation colour, and
niazimicin has no scored row at all because its one biological study dosed an extract
fraction rather than the purified compound. Panel~(b) is deliberately sparse. Only 25
compound-by-gene-set overlaps were recorded (Table~\ref{tab:c4}), 23 for shortlisted
compounds and two for benzyl isothiocyanate, which is not assessable and is not shown. They
are small: 14 involve a single gene and 9 involve two, and the largest are moringin with the
p53 pathway (three genes: BAX, CDKN1A, TP53) and quercetin with the p53 pathway (five genes:
BAX, CDKN1A, DDIT3, TGFB1, TP53). These overlaps are the thinnest kind of link in the
project, and they mark where a hypothesis could be formed, not where an effect has been
shown.

\section{What a tier means, and what it does not}\label{sec:tier-meaning}

A tier records \textbf{the state of the published evidence about a compound}, assembled under
rules fixed before any score was computed (Table~\ref{tab:tier-meaning}). It is not a measure
of radioprotection. Tier~A means \enquote{best documented}, not \enquote{most likely to
work}. Four of the five tier-A compounds are ubiquitous plant metabolites whose radiation
evidence comes entirely from outside \textit{Moringa} research, and the direction of that
evidence differs between them: quercetin and kaempferol have \emph{radiosensitisation}
evidence in tumour cells as well as normal-tissue protection, whereas $\beta$-sitosterol and
rutin have normal-tissue protection only. Radioprotection and radiosensitisation are
opposite effects and are never merged (Table~\ref{tab:c2}).

\begin{table}[tbp]
  \caption{What a compound tier does and does not say.}\label{tab:tier-meaning}
  \small
  \begin{tabular}{@{}L{0.47\linewidth}L{0.47\linewidth}@{}}
    \toprule
    \thead{A tier does say} & \thead{A tier does not say} \\
    \midrule
    Primary studies of the purified compound exist, with the stated independence
      & That the compound protects or sensitises anything \\
    The compound is documented in \textit{Moringa} seed or leaf at the stated
      identification strength & That it was present in the extracts of the original study \\
    Published work connects it to a radiation-response process
      & That it modifies clonogenic survival, the endpoint modelled here \\
    A follow-up experiment can be specified from the record
      & That such an experiment would succeed \\
    \bottomrule
  \end{tabular}
\end{table}

\section{Realistic first-stage experimental priorities}\label{sec:first-stage}

The ranking answers the question \enquote{what is best documented?}. A first experiment has to
answer a narrower question, so the two orderings differ --- without any score being changed.
Three questions can be separated.

\paragraph{Is any plausibly active compound present in the extracts that were tested?}
This is the project's largest gap and the cheapest to close. The priority candidates are
\textbf{niazimicin} (tier~C, the only compound with a confirmed structural identification in
both the seed and the leaf dichloromethane fractions) and \textbf{moringin} (tier~B, the
signature \textit{Moringa} isothiocyanate). Neither has any ionising-radiation study, so
neither can be offered as a protection candidate; they are priorities for
\textbf{analytical characterisation of the original extracts}, followed by a first radiation
screen if they are found.

\paragraph{Does a compound with existing radiation evidence act on the endpoint modelled
here?} \textbf{$\alpha$-Tocopherol} is the strongest first-stage candidate: it is the only
tier-A compound with a full identity score, and the only compound in the evidence base whose
normal-tissue protection evidence includes a \emph{clonogenic} endpoint (rat intestinal crypt
survival), alongside a randomised trial in radiation mucositis. \textbf{$\beta$-Sitosterol} follows,
with normal-tissue protection reported by two independent groups. Both must be tested in
matched normal and tumour cells, because neither has a tumour arm.

\paragraph{Which candidates should not lead?} \textbf{Chlorogenic acid} should not be a
first-stage protection candidate despite its score: the tumour-protection finding makes it a
cautionary compound in a radiotherapy context. \textbf{Astragalin} is capped for a related
reason, its only DNA-damage-relevant evidence being pro-genotoxic. \textbf{Oleic acid} has a
direct clonogenic test showing \emph{no} effect on radiosensitivity, so it is informative but
not a candidate. Quercetin, kaempferol and rutin are scientifically interesting but
second-stage: their \textit{Moringa} identity rests on tentative dereplication, and the
aglycone--glycoside pharmacokinetic question would have to be settled first.

\begin{keybox}[title={A defensible first-stage programme.}]
Characterise the original extracts (niazimicin, moringin), and test $\alpha$-tocopherol and
$\beta$-sitosterol on clonogenic survival in matched normal and tumour cells over a dose
range. This uses the tiers as summaries of evidence, not as predictions, and leaves every
score exactly as computed.
\end{keybox}

The scientific shape of this module's result is therefore a \textbf{mismatch}, and it is its
most useful output: the compound with the best claim to being present in the tested extracts
has no radiation biology, and the compounds with radiation biology have no confirmed presence
in those extracts. Neither gap can be closed by further literature work; both can be closed by
the experiments proposed in Chapter~\ref{ch:recommendations}.

\chapter{Discussion}\label{ch:discussion}

\section{What the modest improvement means}

Expression features reduced the error of a tissue-only baseline by 2.85\%, consistently in
direction (nine of ten repeats) but small in size and not uniform across tissues, with
Pancreas moving in the opposite direction by nearly three times the pooled gain. The
defensible interpretation is that \textbf{baseline Hallmark pathway activity carries a small
amount of information about intrinsic radiosensitivity beyond tissue of origin --- on this
endpoint, in this cohort and under this model family}. This is not a demonstration that
expression predicts radiosensitivity usefully, and it is not a validated effect: no interval,
no $p$-value, no negative control and no external cohort stand behind it.

The robustness analyses sharpen this reading rather than soften it. The advantage survives a
change of model family only in diminished form, since the random forest gains less (0.00275
against 0.00383). On transformed outcome scales it keeps its direction but nearly disappears
on the neglog scale; a result whose size depends so strongly on the choice of transformation
is small enough to be shaped by it. The one encouraging result, the leave-one-tissue-out
$R^2$ of $+0.031$ for expression alone, comes from an easier comparison on a different design,
and it reverses for two of the nine well-supported tissues (Pancreas $-0.0142$,
Esophagus/Stomach $-0.0030$). Pancreas is unfavourable in both the primary and the
leave-one-tissue-out analyses, so the tissue-level picture, and not only the pooled gain,
must be part of any statement about generalisation.

Three further cautions apply. The comparison is nested (M2 contains M1 as an offset), so the
errors of the two models are strongly correlated and the paired difference, while the right
statistic, is a narrow one. Repeated cross-validation re-splits one fixed set of components
and cannot estimate generalisation to new cohorts. And the baseline itself is weak: M1 is
marginally worse than the global mean on MAE, which indicates that tissue of origin, as
encoded here, contributes little once shrinkage is taken into account.

\section{Why out-of-sample $R^2$ remains low}

Figure~\ref{fig:prediction} shows predictions compressed near the cohort mean against
observations that span the whole interval. Several causes act together and cannot be
separated with these data. The \textbf{label} is a literature median that aggregates
laboratories, dose rates and protocols, and the field's own reporting audit found a third of
papers omitting replicates or dose rate \parencite{oike2020}; CLARIFi supplies one median per
line, so per-line assay variability cannot even be estimated. The \textbf{sample} is small
(237 independent components) and tissue-imbalanced. The \textbf{features} are 50 coarse
pathway summaries of baseline expression, which may simply not encode the determinants of
radiation response; intrinsic radiosensitivity may depend on dynamic, post-irradiation
responses that a baseline profile cannot capture. And the \textbf{design is strict}: grouping
by relatedness and confining all preprocessing to training folds remove the two most common
sources of inflated performance.

This near-zero $R^2$ closely reproduces RadSigBench, which reported mean out-of-sample $R^2$
of 0.01 and 0.00 for seven published signatures and found them no better than random gene
sets \parencite{oconnor2022}. \textbf{Agreement with the strongest prior negative result is a
substantive finding, not a failure}: it indicates that the pipeline measures what it claims
to measure, and it is evidence against the more optimistic claims of the signature
literature.

\section{The value and uncertainty of expression-derived pathways}

The pathway ranking is best read as a description of a weak fitted model. Its magnitudes are
of the same order as the whole predictive gain, its rank stability across fits is modest
(mean Kendall $\tau$ 0.225, mean top-10 Jaccard 0.391), and correlated gene sets share or
exchange importance. The appearance of HALLMARK\_P53\_PATHWAY at the top is biologically
\emph{plausible}, since p53 governs the damage response, cell-cycle arrest and apoptosis, and
its predominantly negative coefficient is consistent with higher p53-pathway activity
accompanying lower survival. Plausibility is, however, precisely the trap: without a
demonstrated prediction gain and without negative controls, a plausible ranking cannot be
distinguished from a coincidence, and the absence of DNA repair, G2M checkpoint and hypoxia
from the shortlist shows that the ranking is not simply recovering textbook radiobiology. The
ranking is therefore reported as capped, exploratory and hypothesis-generating, and no
mechanistic conclusion is drawn from it.

\section{How the original \textit{Moringa} findings motivate further work}

The original study's lasting contribution is a \emph{question} rather than a measurement this
project could use: the post-irradiation proliferative response to a plant extract differed
between cell types, most clearly for the leaf extract (PEF 1.26 in V79 against 0.90 and 0.99
in the two carcinoma lines). The computational work formalises the general version of that
question on a valid endpoint and with proper grouping, and finds the molecular signal weak.
This does not contradict the original observation, because the endpoints differ: a 72-hour
biomass ratio measured with extract present is not intrinsic clonogenic survival. The
original observation also remains confounded, since the only apparently normal line is the
only non-human one. What the project adds is a defensible framework in which the question can
be asked properly, together with an explicit statement of the experiments needed to answer it.

\section{What compound prioritisation adds, and why the candidates remain hypotheses}

The compound module turns an open-ended phytochemical literature into an auditable evidence
trail: compound $\rightarrow$ documented presence in seed or leaf with a named solvent
$\rightarrow$ primary experimental association with a radiation-response process
$\rightarrow$ the genes actually measured. Its practical value lies as much in what it
\emph{excludes} as in what it ranks. It removes benzyl isothiocyanate (a root, not a seed or
leaf, constituent), gallic and caffeic acid (no primary \textit{Moringa} report), a UVB
literature that is not radiation evidence, carrier-confounded nanoparticle studies and a
retracted paper; it places the free flavonol aglycones in the weakest identification band
rather than excluding them, as their tentative dereplication warrants; and it flags the
tumour-protection liability of chlorogenic acid, which a naive \enquote{antioxidants are
radioprotectors} reading would have turned into a selling point.

The candidates remain hypotheses for specific reasons. No compound has been shown to be
present in the extracts that were tested, because those extracts were never analysed.
Moringin, the most \textit{Moringa}-specific candidate, has no radiation study at all. The
compounds that do have radiation evidence were mostly tested at concentrations or doses whose
relevance to an extract is unknown, several at levels well above plausible dietary exposure,
and many abstracts omit the radiation dose or the concentration. An antioxidant is not
thereby a radioprotector; a compound that reduces tumour-cell viability does not thereby
protect normal tissue; and a high priority score reflects the state of the published
evidence, not a predicted experimental outcome.

The link to the machine-learning results is deliberately kept weak, because the evidence for
it is weak. With the gene-set ranking capped at tier~C, the compound prioritisation rests on
external biological evidence, and any overlap between a compound's measured genes and a
ranked gene set is reported as a hypothesis. \textbf{Nothing in this report claims that a
prioritised pathway and a prioritised compound together imply an effect on clonogenic SF2.}

\section{Separating supported results from plausible explanations}

Table~\ref{tab:supported} separates what the saved outputs support from explanations that are
plausible but not established. The left-hand column is the evidential content of this report;
the right-hand column is the agenda for future work.

\begin{table}[tbp]
  \caption{Results supported by the saved outputs, separated from plausible explanations
    that are not established.}\label{tab:supported}
  \small
  \begin{tabular}{@{}L{0.47\linewidth}L{0.47\linewidth}@{}}
    \toprule
    \thead{Supported by the outputs} & \thead{Plausible, not established} \\
    \midrule
    $\Delta$MAE $=0.00383$ (2.85\%); M2 better in nine of ten repeats
      & That expression features would improve prediction in an independent cohort \\
    Out-of-sample $R^2$ near zero for all five models
      & That richer features or more data would raise it materially \\
    M2 lower error in eight of nine well-supported tissues, worse in Pancreas
      & That tissue-specific differences reflect tissue-specific biology \\
    Expression-only transfer to unseen tissues: $R^2=+0.031$, negative for Pancreas and
      Esophagus/Stomach & That expression generalises to tissues absent from training \\
    Advantage keeps its direction on three outcome scales but nearly vanishes on one
      & That the raw-scale magnitude is the \enquote{true} size of the effect \\
    Seven gene sets meet the tier-C criteria; the p53 pathway ranks first
      & That the p53 pathway influences radiosensitivity in these lines \\
    Importance rank stability is modest across fits
      & That a larger sample would stabilise the ranking \\
    Niazimicin is confirmed in seed and leaf DCM fractions
      & That it was present in the original study's extracts \\
    Moringin has no published ionising-radiation study
      & That it would behave like benzyl isothiocyanate under radiation \\
    Chlorogenic acid protected hepatocellular carcinoma cells from radiation in one study
      & That it would protect human tumours during radiotherapy \\
    \bottomrule
  \end{tabular}
\end{table}

\chapter{Implementation, challenges and deliverables}\label{ch:implementation}

\section{Challenges and how they were addressed}

\paragraph{An incomplete source manuscript.} Printed page~3 and pages 7--12 of the original
manuscript are missing, removing the extraction and culture methods, the end of the
discussion, the conclusions and the references; authors, venue and DOI are absent, and no
published version was found. In the manuscript's own figures the 0~Gy and 2~Gy series cannot
be distinguished in the greyscale scan. The response was a labelled page-by-page extraction
that marks every absent item as not stated, with no reconstruction and no use of point values
read from the figures.

\paragraph{Endpoint confusion as a standing risk.} The central scientific hazard of the
project is to treat proliferation, viability and clonogenic survival as interchangeable.
Endpoints are therefore separated in the data model, the documentation and the figure rules;
a dedicated comparison of PEF with SF2 is recorded; and a project rule forbids placing
different endpoints on one axis.

\paragraph{Dataset harmonisation and licensing.} Four resources with different identifier
namespaces, release cycles and licences had to be joined, and the outcome dataset is licensed
for non-commercial use only. Each source has a manifest with checksums and access dates, the
loaders verify checksums, and licences are recorded explicitly. The identity of the Hallmark
gene-set file with the official release remains to be confirmed.

\paragraph{Cell-line identity and relatedness.} Names are ambiguous, and relatedness is the
most dangerous leakage route in this design. Matching used stable identifiers only;
Cellosaurus hierarchy and origin links merged 241 records into 237 components; problem flags
on relatives were recorded and surfaced in the triage; and the identifier, component and
split code was audited for leakage.

\paragraph{A small sample against many features.} With 237 independent components only modest
model complexity is supportable. Simple scikit-learn baselines were used before anything more
complex, within a prespecified ladder of five models with nested tuning, and the random forest
was confined to an explicitly post-hoc exploratory configuration.

\paragraph{Computational cost.} The validation plan required a coverage simulation to choose an
interval method before any inference. The first pilot took about 24--25 minutes per replicate
on one core, which made $9\times500$ replicates impractical within the available time. A
static review found redundant work in the inner tuning loop: the gene-set scorer, the score
scaler and the validation scores were being refitted for each of the 36 grid points in each of
five inner folds, although they depend only on the fold rows, and the baseline ridge was
refitted per grid point although it depends only on its own penalty. A caching rewrite was
designed to be bit-identical (same operations in the same order), with the original loop kept
as the reference implementation for every model to which the cache does not apply. Bounded
multiprocessing was added (spawn context, atomic write-once checkpoints written by the parent
process only, one BLAS thread per worker), sized to the investigator's machine (six physical
cores, twelve threads, about 31~GiB of memory). Even so, the full simulation could not be
completed before the deadline.

\section{The protocol deviation}\label{sec:deviation}

The deviation was written into the decision log \textbf{before} the primary run produced a
single result, which is what makes it a deviation rather than a rationalisation. Its terms
are: the inference method is set to \enquote{not assessable}; no interval or $p$-value is
written; the number of repeats is fixed at ten; the negative controls are left at zero draws;
and, as a consequence, gene-set attribution can never be permitted and the pathway tiers stay
capped. The deviation is reflected in the results, the figures and the limitations, and the
simulation infrastructure remains in the repository ready to run (Table~\ref{tab:unexecuted}).

\begin{table}[tbp]
  \caption{Analyses implemented and tested but not executed.}\label{tab:unexecuted}
  \small
  \begin{tabular}{@{}L{0.24\linewidth}L{0.33\linewidth}L{0.36\linewidth}@{}}
    \toprule
    \thead{Analysis} & \thead{Command} & \thead{Status} \\
    \midrule
    Negative controls, 1,000 draws of each null
      & \path{scripts/run_model.py} with \path{configs/experiment_controls.yaml}
      & Configuration ready, cost measured at about 10--20~h; \textbf{not executed}. Tiers A
        and B stay closed and no gene-set $p$-value exists until it is run. \\
    Coverage simulation (interval method) & \path{scripts/run_interval_simulation.py}
      & Infrastructure ready; \textbf{deferred} by the recorded protocol deviation. \\
    \bottomrule
  \end{tabular}
\end{table}

\paragraph{A presentation error found and corrected.} The first version of the within-tissue
Spearman statistic was strongly negative for the tissue-only models. Rather than report a
striking number, the code was inspected, the leave-fold-out mean artefact identified, and a
corrected table written that marks the statistic undefined where a model's prediction is
constant within a tissue (Section~\ref{sec:correction}; Table~\ref{tab:s4}). The first version
was kept unchanged for audit and is not interpreted --- a clear example of the write-once
discipline doing its job.

\paragraph{Uncertainty in the correspondence between compounds and extracts.} The original
extracts were never chemically characterised, and the published dichloromethane-extract
literature is thin, partly inaccessible and partly based on library matches. The response was
a three-level extraction category per compound, kept outside the score; verbatim solvent
wording stored for every report; the exclusion of speculation as evidence (one paper's
\enquote{possibly yes} for chlorogenic acid in a leaf \ce{CH2Cl2} extract); and a note of a
citation error in a 2026 paper, in which a statement about an Indian DCM extract cites an
unrelated soil-chemistry reference.

\paragraph{Abstract-level evidence extraction.} Many full texts were paywalled or could not be
retrieved, so radiation doses, concentrations and even cell-line identities are missing from
several records. Every such gap is recorded as \enquote{not stated} together with the reading
depth, and the gaps are listed as items to close before publication rather than filled by
inference.

\section{Deliverables}\label{sec:deliverables}

The project delivered the following outputs.
\begin{enumerate}
  \item A labelled extraction of the original manuscript, with every absent item marked and
    checked against the page images.
  \item A provenance-tracked analysis dataset of 241 records and 237 independent components
    from 20 tissues, with checksum-verified inputs and a recorded quality-control pass.
  \item The primary benchmark: 250/250 model-fold tasks for five prespecified models in ten
    repeats of grouped, nested five-fold cross-validation, with saved out-of-fold predictions,
    metrics, per-repeat differences and a run record.
  \item The explainability analysis: held-out permutation importance for all 50 fits of M2,
    coefficient sign consistency, gene-set score correlations and rank-stability statistics.
  \item Pathway and cell-line triage under enforced prespecified gates: seven tier-C gene
    sets and no tier A or B; 12, 9 and 27 cell lines in tiers A, B and C.
  \item Publication tables derived only from saved outputs, including the corrected
    within-tissue ranking analysis with the superseded version retained for audit.
  \item A curated, fully referenced compound evidence base (constituent inventory and primary
    radiation-biology evidence), with every DOI verified against Europe PMC or Crossref and
    every identifier against PubChem.
  \item The prespecified compound-prioritisation protocol, its tested implementation and the
    resulting ranking: 18 compounds ranked, 5 in tier~A, with the liability cap and the
    not-assessable rule both binding.
  \item Three completed robustness analyses --- the M3 random forest (200/200 tasks),
    leave-one-tissue-out (18/18) and outcome-scale sensitivity (400/400) --- each reported as
    it came out, including where it weakened the primary reading, and M3 explainability kept
    separate from the primary analysis.
  \item Eleven report figures as 600~dpi PNG and vector PDF with provenance records, and this
    report.
\end{enumerate}
All model fitting, prediction, explainability and analysis were executed by the
investigator; the automated assistant prepared code, configuration, documentation and this
report, and ran no scientific computation.

\chapter{Limitations}\label{ch:limitations}

\begin{enumerate}
  \item \textbf{Predictive performance is weak.} Out-of-sample $R^2$ lies between $-0.018$ and
    $+0.013$ for every model. No model in this study predicts clonogenic radiosensitivity
    usefully, and the pipeline should not be presented as a working predictor.
  \item \textbf{The improvement is small and was not formally tested.} $\Delta\mathrm{MAE}=
    0.00383$ mSF2 units (2.85\%). No interval, $p$-value or formal superiority assessment
    exists, because the coverage simulation that would have validated an interval method was
    deferred.
  \item \textbf{No negative controls were run.} Zero permuted-outcome and zero random-gene-set
    draws were completed, so the gain cannot be compared with a null reference, gene-set
    attribution is not permitted and the stability gate cannot be evaluated.
  \item \textbf{There is no external validation.} The ten repeats re-split the same 237
    components; repeated cross-validation measures split-to-split variability, not
    generalisation.
  \item \textbf{The improvement is not uniform, and transfer to unseen tissues is
    conditional.} In Pancreas, M2 is worse than M1 by 0.0110, nearly three times the pooled
    gain; in the leave-one-tissue-out analysis the expression-only model lost to the global
    mean when Pancreas ($-0.0142$) or Esophagus/Stomach ($-0.0030$) was held out.
  \item \textbf{The size of the advantage depends on the outcome scale}, and magnitudes on
    the neglog and logit scales are in transformed units that cannot be compared directly
    with the raw-scale estimate.
  \item \textbf{One outcome source with aggregated labels.} CLARIFi mSF2 is a literature
    median across laboratories, dose rates and protocols, with one value per line, so per-line
    assay variability cannot be estimated; the dataset is also licensed for non-commercial use
    only.
  \item \textbf{A small, imbalanced sample.} There are 237 components; eleven of the 20
    tissues have fewer than ten, and one has a single component.
  \item \textbf{Pathway rankings are exploratory and only moderately stable.} Mean Kendall
    $\tau$ is 0.225 and mean top-10 Jaccard 0.391; importance magnitudes are the size of the
    whole predictive gain; correlated sets share importance; no tier exceeds C.
  \item \textbf{Cell-line discrepancies are dominated by shrinkage.} Large residuals occur
    mostly at extreme observed values and are review signals, not biological findings;
    measured lines are not independent validation samples.
  \item \textbf{No compound has been shown to be present in the original extracts.} The
    extracts were never analysed; most compounds are documented with other solvents; and
    several DCM identifications are spectral-library matches from inaccessible full texts,
    one in a journal indexed in neither PubMed nor Europe PMC.
  \item \textbf{The most \textit{Moringa}-specific candidate has no radiation evidence.} No
    primary study has combined purified moringin or glucomoringin with ionising radiation.
  \item \textbf{Compound evidence is largely abstract-level.} Radiation doses, concentrations
    and some cell-line identities are missing from several records and must be recovered from
    the full texts; the glucomoringin identifier remains ambiguous; the erratum to
    \textcite{abedpour2022} has not been read; and because the journal that published
    \textcite{yin2022} had a large contemporaneous retraction wave, that article's status must
    be re-checked (no retraction notice was found on 22 October 2026).
  \item \textbf{Group independence is lower than paper counts suggest.} Shared senior authors
    and a shared compound supplier link several studies; the scoring counts research groups
    once, but the effective number of independent sources remains small.
  \item \textbf{The within-tissue ranking statistic is unusable for the simpler models.} It is
    undefined where predictions are constant within a tissue and biased downward elsewhere.
  \item \textbf{No causal or mechanistic claim is supported.} Association, explanation of a
    fitted model, causation and experimentally demonstrated protection are four different
    things, and only the first two are in evidence here.
\end{enumerate}

\chapter{Recommendations for experimental validation}\label{ch:recommendations}

The following priorities are proposals; none has been carried out. They are ordered so that
each removes an uncertainty on which the next depends.

\paragraph{Priority 1 --- close the chemistry gap.} Characterise the \emph{original}
dichloromethane seed and leaf extracts by LC-MS/MS and GC-MS against authentic standards,
using the published validated UPLC-MS/MS method for niazimicin \parencite{abdelsayed2021jfda}
as a positive control. Until this is done, no compound-level interpretation of the original
result is possible. Yields and concentrations should be reported as measured, never inferred.

\paragraph{Priority 2 --- resolve identities.} Confirm the structures and public identifiers of
the shortlisted compounds, settle the glucomoringin stereochemistry, obtain the paywalled
full texts and recover the missing radiation doses and concentrations. Re-verify the
retraction status of the flagged journal and read the outstanding erratum.

\paragraph{Priority 3 --- complete the deferred statistical validation.} Run the coverage
simulation to choose and justify an interval method, re-assess the primary comparison with
the adopted method, and run at least 1,000 draws of each negative control (permuted outcome
and random gene sets). Only then can superiority be assessed, gene-set attribution be
permitted or any pathway exceed tier~C. The infrastructure for all of this is in place.

\paragraph{Priority 4 --- test purified compounds against whole extracts on the right
endpoint.} Use multi-dose \textbf{clonogenic survival} assays (not 72-hour biomass) with and
without compound or extract, yielding linear-quadratic parameters,
\begin{equation}
  \mathrm{SF}(D) = \exp\bigl(-\alpha D-\beta D^{2}\bigr),
  \label{eq:lq}
\end{equation}
where $\alpha$ and $\beta$ here denote the linear and quadratic radiobiological coefficients,
and a dose-modifying factor for each treatment. Include vehicle-only and cell-free
interference controls, which the original study did not describe, and independent biological
repeats on separate days and passages, since the original replicates were technical wells and
dose was confounded with plate.

\paragraph{Priority 5 --- make the normal-versus-tumour comparison interpretable.} Use
species-matched and tissue-matched normal and tumour pairs assayed identically, so that
species, tissue and malignancy are not confounded as they were when one hamster fibroblast
line was compared with two human carcinoma lines. A selectivity claim requires treated normal
\emph{and} treated tumour measurements on the same endpoint. The best published model of how
such a claim should be worded is the quercetin study of \textcite{horton2013}, which reports
normal-tissue anti-fibrotic protection with \emph{no} effect on fibroblast or murine tumour
radioresponse.

\paragraph{Priority 6 --- test the prioritised pathways directly, if Priority 3 succeeds.}
Perturb the implicated processes (for example p53-pathway activity) in a small panel of cell
lines spanning the predicted range, and measure clonogenic survival. This is hypothesis
testing of a model-derived conjecture, not confirmation.

\paragraph{Priority 7 --- strengthen the benchmark.} Seek a second, independent
radiosensitivity source for cross-assay evaluation, reported as heterogeneity rather than
pooled. Consider prospective prediction on unassayed lines under a panel-selection rule
written into the protocol \emph{before} predictions are viewed, reported as a designed
validation panel rather than an unbiased external cohort.

\chapter{Conclusions}\label{ch:conclusion}

This project set out to connect an unpublished \textit{Moringa oleifera} irradiation study to
a defensible computational investigation of cellular radiosensitivity, and to use the result
to identify compounds worth testing.

\paragraph{What was established.} A reproducible, provenance-tracked pipeline links public
clonogenic radiosensitivity data to baseline expression for 237 independent biological
components, with cell-line relatedness handled correctly and all preprocessing confined to
training folds. Within that design, pathway-level expression features reduce the error of a
tissue-only baseline by 0.00383 mSF2 units (2.85\%), consistently in direction across nine of
ten cross-validation repeats but small, non-uniform across tissues and unvalidated by any
formal test. Overall predictive power is weak: out-of-sample $R^2$ is near zero for every
model, closely reproducing the strongest published negative benchmark in the field.

Three further analyses of the same cohort tested how robust that advantage is. A random
forest did worse than the linear offset model; on transformed outcome scales the advantage
kept its direction but nearly vanished on one; and an expression-only model transferred to
unseen tissues with the only clearly positive out-of-sample $R^2$ of the project, but lost to
the global mean when Pancreas or Esophagus/Stomach was held out.

The compound module produced a curated, fully referenced evidence base on \textit{Moringa}
constituents in which the extraction solvent is tracked explicitly, and a shortlist of 18
candidates with 5 in tier~A, in which the prespecified liability cap and the not-assessable
rule both visibly bound. It established that niazimicin is the only constituent confirmed in a
dichloromethane fraction of seed and leaf, and that moringin, the isothiocyanate most
characteristic of the plant, has no published ionising-radiation study at all.

\paragraph{What remains hypothetical.} Everything downstream of the prediction. The pathway
ranking is capped at an exploratory tier because no prediction gain was formally shown and no
negative controls were run. The cell-line discrepancies are review signals dominated by
shrinkage. No compound has been shown to be present in the extracts that were tested, no
compound has been shown in this work to modify clonogenic survival, and no claim of selective
radioprotection is made or supported. The candidate tiers describe the state of the published
evidence, not expected experimental outcomes.

\paragraph{Contribution.} The project's contribution is methodological and evidential rather
than predictive: a transparent framework; an honest, negative-leaning benchmark that agrees
with the best prior work; the correction of its own earlier artefact; an explicit record of a
deadline-driven protocol deviation; and a traceable path from the original radiobiology to a
small set of experiments that would settle the question. Its most valuable single output may
be the mismatch revealed by the compound curation --- the compound present in the right
solvent has no radiation data, and the compounds with radiation data were never found in that
solvent --- because it turns a diffuse phytochemical question into two specific and feasible
experiments: characterise the original extracts, and test the best-documented candidates on
clonogenic survival in matched normal and tumour cells.


\printbibliography[heading=bibintoc,title={References}]

% --------------------------------------------------------------------- appendices
\reportappendix
% Appendices. Supplementary tables and figures are numbered S1, S2, ... .

% =====================================================================================
\section{Supplementary results of the primary benchmark}\label{app:primary}

\begin{table}[htbp]
  \caption[$\Delta$MAE per cross-validation repeat]{Paired, component-weighted $\Delta$MAE
    (mSF2 units) in each cross-validation repeat. Repeats re-split the same lines, so their
    spread is split-to-split variability, not sampling uncertainty. Source:
    \path{results/tables/primary_v1/publication_v1/table_s1_repeat_delta_mae.csv}.}
  \label{tab:s1}
  \small
  \begin{tabular}{@{}l S[table-format=-1.6] S[table-format=1.6]@{}}
    \toprule
    \thead{Repeat} & {\thead{M2 vs M1 (primary)}} & {\thead{M2P vs M1P (sensitivity)}} \\
    \midrule
    1 & 0.006258 & 0.004976 \\
    2 & 0.003042 & 0.003706 \\
    3 & 0.005528 & 0.004016 \\
    4 & 0.002201 & 0.002335 \\
    5 & 0.002619 & 0.002303 \\
    6 & 0.005018 & 0.004040 \\
    7 & -0.002404 & 0.000379 \\
    8 & 0.007638 & 0.007469 \\
    9 & 0.005673 & 0.005690 \\
    10 & 0.002742 & 0.003793 \\
    \midrule
    Mean (pooled estimate) & 0.003832 & 0.003871 \\
    Repeats with $\Delta$MAE $>0$ & {9} & {10} \\
    \bottomrule
  \end{tabular}
\end{table}

\begin{table}[htbp]
  \caption[Within-tissue Spearman correlation (corrected version)]{Within-tissue Spearman
    correlation between observed and pooled out-of-fold predicted mSF2, corrected version,
    for the nine tissues with at least ten components (the eleven smaller tissues are not
    reported). \emph{Undefined}: the prediction is constant within this tissue in every
    fitted model, so there is no within-tissue ranking. Defined values are biased downward,
    because pooled out-of-fold predictions contain each fold's training tissue mean, which is
    lower when the held-out line's own value is high; they are descriptive only. Source:
    \path{results/tables/primary_v1/publication_v1/within_tissue_spearman_v2.csv}.}
  \label{tab:s4}
  \small
  \begin{tabular}{@{}l S[table-format=2.0] c c S[table-format=-1.3] S[table-format=-1.3]
      S[table-format=-1.3]@{}}
    \toprule
    \thead{Tissue} & {\thead{Components}} & \thead{M0} & \thead{M1} & {\thead{M1P}}
      & {\thead{M2}} & {\thead{M2P}} \\
    \midrule
    Bladder/Urinary Tract & 11 & \textit{undefined} & \textit{undefined} & -0.227 & 0.182 & 0.182 \\
    Bowel & 21 & \textit{undefined} & \textit{undefined} & -0.263 & -0.035 & -0.047 \\
    Breast & 17 & \textit{undefined} & \textit{undefined} & -0.507 & 0.252 & 0.223 \\
    CNS/Brain & 23 & \textit{undefined} & \textit{undefined} & -0.046 & 0.333 & 0.333 \\
    Esophagus/Stomach & 23 & \textit{undefined} & \textit{undefined} & -0.294 & -0.179 & -0.211 \\
    Head and Neck & 18 & \textit{undefined} & \textit{undefined} & 0.252 & 0.153 & 0.159 \\
    Lung & 46 & \textit{undefined} & \textit{undefined} & -0.118 & -0.038 & -0.032 \\
    Ovary/Fallopian Tube & 10 & \textit{undefined} & \textit{undefined} & -0.709 & 0.200 & 0.200 \\
    Pancreas & 16 & \textit{undefined} & \textit{undefined} & -0.156 & -0.262 & -0.274 \\
    \bottomrule
  \end{tabular}
\end{table}

\begin{figure}[htbp]
  \centering
  \includegraphics[width=\linewidth,height=0.42\textheight,keepaspectratio]{fig_s01_prediction_variability}
  \caption[Stability of out-of-fold predictions across repeats]{Stability of each cell line's
    out-of-fold prediction across the ten cross-validation repeats, by model: the
    distribution over records of the standard deviation of the prediction. Violins show the
    kernel density over all 241 records with the quartiles marked. The quantity is
    split-to-split variability of the prediction, not a prediction or confidence interval,
    and it widens with the number of repeats. The expression models vary more between splits
    than the tissue-only models, whose prediction is a shrunken tissue mean.}
  \label{fig:s-variability}
\end{figure}

% =====================================================================================
\section{Supplementary results of the explainability analysis}\label{app:explain}

\begin{figure}[htbp]
  \centering
  \includegraphics[width=0.9\linewidth,height=0.5\textheight,keepaspectratio]{fig_s02_gene_set_correlation}
  \caption[Correlation between the scores of the tier-C gene sets and related sets]{Spearman
    correlation between the Hallmark gene-set scores of the tier-C sets and of every set
    flagged as sharing signal with them (absolute correlation at least 0.80), computed over
    the eligible cell lines from the expression features alone. Rows and columns follow the
    triage rank order, with no clustering, so no grouping should be read as a biological
    relationship. Correlated scores can share or exchange permutation importance, so
    neighbouring sets in Figure~\ref{fig:importance} are not independent findings.}
  \label{fig:s-correlation}
\end{figure}

\begingroup\footnotesize\setlength{\tabcolsep}{3.5pt}
\begin{longtable}{@{}r l S[table-format=-1.3] S[table-format=1.3] S[table-format=1.2]
    S[table-format=1.2] l S[table-format=1.2] l L{0.20\linewidth}@{}}
  \caption[Exploratory Hallmark gene-set ranking]{Exploratory Hallmark gene-set ranking with
    prespecified triage tiers (prefix \texttt{HALLMARK\_} omitted). Importance is the
    increase in component-weighted held-out MAE when a score is permuted, averaged over the
    50 fits of M2, in units of $10^{-3}$ mSF2. Shares are proportions of cross-validation fits,
    not probabilities. Tiers A and B are closed (no formally shown gain; zero null draws);
    negative-control evidence is not available for any set. A gene-set-level association
    inside a predictive model is not a gene-level finding, a validated biomarker, a causal
    mechanism or a radiation-response pathway. Source:
    \path{results/tables/primary_v1/publication_triage_v1/table_s2_gene_set_ranking.csv}.}
  \label{tab:s2}\\
  \toprule
  \thead{Rank} & \thead{Gene set} & {\thead{Mean imp.}} & {\thead{SD}}
    & {\thead{Share $>0$}} & {\thead{Sign cons.}} & \thead{Sign} & {\thead{Share zero}}
    & \thead{Tier} & \thead{Score correlated with} \\
  \midrule
  \endfirsthead
  \multicolumn{10}{@{}l}{\tcont{Table~\thetable\ (continued)}}\\
  \toprule
  \thead{Rank} & \thead{Gene set} & {\thead{Mean imp.}} & {\thead{SD}}
    & {\thead{Share $>0$}} & {\thead{Sign cons.}} & \thead{Sign} & {\thead{Share zero}}
    & \thead{Tier} & \thead{Score correlated with} \\
  \midrule
  \endhead
  \midrule
  \multicolumn{10}{@{}r}{\tcont{Continued on next page}}\\
  \endfoot
  \bottomrule
  \endlastfoot
  1 & P53\_PATHWAY & 3.428 & 3.017 & 0.90 & 0.92 & negative & 0.08 & C & \na \\
  2 & COMPLEMENT & 1.181 & 2.479 & 0.68 & 0.92 & positive & 0.08 & C & APOPTOSIS \\
  3 & OXIDATIVE\_PHOSPHORYLATION & 1.070 & 1.639 & 0.74 & 0.86 & negative & 0.14 & C & \na \\
  4 & UV\_RESPONSE\_UP & 1.000 & 2.159 & 0.68 & 0.76 & negative & 0.24 & C & PI3K\_AKT\_MTOR\_SIGNALING \\
  5 & MYC\_TARGETS\_V2 & 0.638 & 1.150 & 0.64 & 0.88 & negative & 0.12 & C & MYC\_TARGETS\_V1 \\
  6 & KRAS\_SIGNALING\_UP & 0.620 & 1.577 & 0.58 & 0.82 & positive & 0.18 & C & \na \\
  7 & UV\_RESPONSE\_DN & 0.551 & 1.558 & 0.60 & 0.82 & positive & 0.18 & C & TGF\_BETA\_SIGNALING \\
  \midrule
  8 & EPITHELIAL\_MESENCHYMAL\_TRANSITION & 0.393 & 1.118 & 0.38 & 0.56 & positive & 0.44 & \na & \na \\
  9 & NOTCH\_SIGNALING & 0.262 & 1.119 & 0.40 & 0.52 & negative & 0.48 & \na & \na \\
  10 & HYPOXIA & 0.077 & 0.514 & 0.10 & 0.04 & negative & 0.96 & \na & UV\_RESPONSE\_DN \\
  11 & INTERFERON\_ALPHA\_RESPONSE & 0.052 & 0.353 & 0.10 & 0.04 & negative & 0.96 & \na & INTERFERON\_GAMMA\_RESPONSE \\
  12 & KRAS\_SIGNALING\_DN & 0.038 & 0.591 & 0.40 & 0.66 & negative & 0.34 & \na & \na \\
  13 & G2M\_CHECKPOINT & 0.011 & 0.075 & 0.08 & 0.02 & positive & 0.98 & \na & E2F\_TARGETS \\
  14 & ADIPOGENESIS & 0.003 & 0.051 & 0.12 & 0.10 & negative & 0.90 & \na & FATTY\_ACID\_METABOLISM \\
  15 & TNFA\_SIGNALING\_VIA\_NFKB & 0.002 & 0.013 & 0.08 & 0.02 & negative & 0.98 & \na & INFLAMMATORY\_RESPONSE \\
  16 & ANDROGEN\_RESPONSE & 0.000 & 0.000 & 0.06 & 0.00 & none & 1.00 & \na & PROTEIN\_SECRETION \\
  17 & APOPTOSIS & 0.000 & 0.000 & 0.06 & 0.00 & none & 1.00 & \na & COMPLEMENT \\
  18 & E2F\_TARGETS & 0.000 & 0.000 & 0.06 & 0.00 & none & 1.00 & \na & G2M\_CHECKPOINT \\
  19 & FATTY\_ACID\_METABOLISM & 0.000 & 0.000 & 0.06 & 0.00 & none & 1.00 & \na & ADIPOGENESIS \\
  20 & GLYCOLYSIS & 0.000 & 0.000 & 0.06 & 0.00 & none & 1.00 & \na & PROTEIN\_SECRETION \\
  21 & HEME\_METABOLISM & 0.000 & 0.000 & 0.06 & 0.00 & none & 1.00 & \na & PROTEIN\_SECRETION \\
  22 & IL2\_STAT5\_SIGNALING & 0.000 & 0.000 & 0.06 & 0.00 & none & 1.00 & \na & APOPTOSIS \\
  23 & IL6\_JAK\_STAT3\_SIGNALING & 0.000 & 0.000 & 0.06 & 0.00 & none & 1.00 & \na & INFLAMMATORY\_RESPONSE \\
  24 & INTERFERON\_GAMMA\_RESPONSE & 0.000 & 0.000 & 0.06 & 0.00 & none & 1.00 & \na & INTERFERON\_ALPHA\_RESPONSE \\
  25 & MITOTIC\_SPINDLE & 0.000 & 0.000 & 0.06 & 0.00 & none & 1.00 & \na & G2M\_CHECKPOINT \\
  26 & MTORC1\_SIGNALING & 0.000 & 0.000 & 0.06 & 0.00 & none & 1.00 & \na & UNFOLDED\_PROTEIN\_RESPONSE \\
  27 & MYC\_TARGETS\_V1 & 0.000 & 0.000 & 0.06 & 0.00 & none & 1.00 & \na & MYC\_TARGETS\_V2 \\
  28 & PEROXISOME & 0.000 & 0.000 & 0.06 & 0.00 & none & 1.00 & \na & BILE\_ACID\_METABOLISM \\
  29 & PI3K\_AKT\_MTOR\_SIGNALING & 0.000 & 0.000 & 0.06 & 0.00 & none & 1.00 & \na & MITOTIC\_SPINDLE \\
  30 & SPERMATOGENESIS & 0.000 & 0.000 & 0.06 & 0.00 & none & 1.00 & \na & \na \\
  31 & TGF\_BETA\_SIGNALING & 0.000 & 0.000 & 0.06 & 0.00 & none & 1.00 & \na & UV\_RESPONSE\_DN \\
  32 & UNFOLDED\_PROTEIN\_RESPONSE & 0.000 & 0.000 & 0.06 & 0.00 & none & 1.00 & \na & MTORC1\_SIGNALING \\
  33 & ESTROGEN\_RESPONSE\_EARLY & -0.002 & 0.024 & 0.08 & 0.04 & negative & 0.96 & \na & ESTROGEN\_RESPONSE\_LATE \\
  34 & CHOLESTEROL\_HOMEOSTASIS & -0.005 & 0.033 & 0.06 & 0.02 & positive & 0.98 & \na & \na \\
  35 & ANGIOGENESIS & -0.008 & 0.124 & 0.08 & 0.04 & positive & 0.94 & \na & COAGULATION \\
  36 & APICAL\_JUNCTION & -0.008 & 0.059 & 0.06 & 0.02 & positive & 0.98 & \na & TGF\_BETA\_SIGNALING \\
  37 & MYOGENESIS & -0.013 & 0.429 & 0.06 & 0.04 & negative & 0.96 & \na & APICAL\_JUNCTION \\
  38 & PANCREAS\_BETA\_CELLS & -0.014 & 0.070 & 0.06 & 0.02 & none & 0.96 & \na & \na \\
  39 & INFLAMMATORY\_RESPONSE & -0.015 & 0.107 & 0.08 & 0.12 & positive & 0.88 & \na & TNFA\_SIGNALING\_VIA\_NFKB \\
  40 & ALLOGRAFT\_REJECTION & -0.026 & 0.184 & 0.06 & 0.02 & positive & 0.98 & \na & APOPTOSIS \\
  41 & WNT\_BETA\_CATENIN\_SIGNALING & -0.028 & 0.198 & 0.12 & 0.20 & negative & 0.80 & \na & \na \\
  42 & ESTROGEN\_RESPONSE\_LATE & -0.036 & 0.202 & 0.06 & 0.04 & positive & 0.96 & \na & ESTROGEN\_RESPONSE\_EARLY \\
  43 & BILE\_ACID\_METABOLISM & -0.037 & 0.264 & 0.06 & 0.02 & positive & 0.98 & \na & PEROXISOME \\
  44 & XENOBIOTIC\_METABOLISM & -0.041 & 0.232 & 0.06 & 0.04 & positive & 0.96 & \na & \na \\
  45 & REACTIVE\_OXYGEN\_SPECIES\_PATHWAY & -0.042 & 0.426 & 0.18 & 0.24 & positive & 0.76 & \na & \na \\
  46 & HEDGEHOG\_SIGNALING & -0.049 & 0.271 & 0.06 & 0.08 & positive & 0.92 & \na & \na \\
  47 & DNA\_REPAIR & -0.054 & 0.379 & 0.06 & 0.02 & negative & 0.98 & \na & G2M\_CHECKPOINT \\
  48 & APICAL\_SURFACE & -0.066 & 0.290 & 0.06 & 0.06 & positive & 0.94 & \na & \na \\
  49 & COAGULATION & -0.082 & 0.461 & 0.08 & 0.08 & positive & 0.92 & \na & COMPLEMENT \\
  50 & PROTEIN\_SECRETION & -0.107 & 0.754 & 0.06 & 0.02 & positive & 0.98 & \na & PI3K\_AKT\_MTOR\_SIGNALING \\
\end{longtable}
\tablenote{Tier: C, exploratory tier~C; \na, not prioritised. Sign: dominant sign of the
  coefficient across fits. Share zero: share of fits in which the set received a zero
  coefficient. Values of 0.000 are numerically zero to three decimals (sets never selected).
  \enquote{Score correlated with}: a set whose score correlates with this one at
  $\lvert\rho\rvert\ge0.80$, so that importance may be shared or split between the two.}
\endgroup

% ----------------------------------------------------------------- cell-line triage
\begin{landscape}
\begingroup\scriptsize\setlength{\tabcolsep}{3pt}
\begin{longtable}{@{}r l l L{0.10\linewidth} S[table-format=1.3] S[table-format=1.3]
    S[table-format=1.3] S[table-format=1.3] S[table-format=1.3] c L{0.13\linewidth}
    L{0.28\linewidth}@{}}
  \caption[Cell-line discrepancy review]{Cell-line discrepancy review: every line in
    discrepancy tiers A, B and C, and MCF-7 (not prioritised), shown for context; the other
    192 lines are not prioritised. Predictions are the repeat-averaged out-of-fold
    predictions (mSF2); disagreement is $\lvert$M1$-$M2$\rvert$; SD is the spread of the M2
    prediction across repeats. Tier thresholds are cohort percentiles of the absolute error
    (Section~\ref{sec:triage-methods}). Measured lines are not independent validation
    samples, and context lines are labelled only; their tiers are not changed. Source:
    \path{results/tables/primary_v1/publication_triage_v1/table_s3_discrepancy_review.csv}.}
  \label{tab:s3}\\
  \toprule
  \thead{Rank} & \thead{DepMap ID} & \thead{Cell line} & \thead{Tissue}
    & {\thead{Observed}} & {\thead{Pred.\ M1}} & {\thead{Pred.\ M2}}
    & {\thead{Disagr.}} & {\thead{SD M2}} & \thead{Tier} & \thead{Triage category}
    & \thead{Technical flags} \\
  \midrule
  \endfirsthead
  \multicolumn{12}{@{}l}{\tcont{Table~\thetable\ (continued)}}\\
  \toprule
  \thead{Rank} & \thead{DepMap ID} & \thead{Cell line} & \thead{Tissue}
    & {\thead{Observed}} & {\thead{Pred.\ M1}} & {\thead{Pred.\ M2}}
    & {\thead{Disagr.}} & {\thead{SD M2}} & \thead{Tier} & \thead{Triage category}
    & \thead{Technical flags} \\
  \midrule
  \endhead
  \midrule
  \multicolumn{12}{@{}r}{\tcont{Continued on next page}}\\
  \endfoot
  \bottomrule
  \endlastfoot
  1 & ACH-000322 & HT-144 & Skin & 0.080 & 0.594 & 0.593 & 0.001 & 0.013 & A & technical review first & small tissue group \\
  2 & ACH-000622 & KNS-42 & CNS/Brain & 0.125 & 0.586 & 0.596 & 0.009 & 0.017 & A & unexplained discrepancy & \na \\
  3 & ACH-000941 & HEC-1-B & Uterus & 0.970 & 0.568 & 0.561 & 0.006 & 0.009 & A & technical review first & related lines in the cohort; small tissue group \\
  4 & ACH-000055 & D283 Med & CNS/Brain & 0.180 & 0.585 & 0.582 & 0.003 & 0.023 & A & unexplained discrepancy & \na \\
  5 & ACH-000900 & NCI-H23 & Lung & 0.190 & 0.593 & 0.583 & 0.010 & 0.012 & A & unexplained discrepancy & \na \\
  6 & ACH-000808 & HuH28 & Biliary Tract & 0.998 & 0.583 & 0.625 & 0.042 & 0.014 & A & technical review first & small tissue group \\
  7 & ACH-000048 & TOV-112D & Ovary/Fallopian Tube & 0.970 & 0.584 & 0.546 & 0.038 & 0.019 & A & unexplained discrepancy & \na \\
  8 & ACH-000736 & SNU-601 & Esophagus/Stomach & 0.190 & 0.582 & 0.560 & 0.022 & 0.010 & A & unexplained discrepancy & \na \\
  9 & ACH-000335 & MSTO-211H & Pleura & 0.970 & 0.591 & 0.579 & 0.012 & 0.013 & A & technical review first & small tissue group \\
  10 & ACH-000934 & MDA-MB-361 & Breast & 0.180 & 0.582 & 0.542 & 0.040 & 0.012 & A & unexplained discrepancy & \na \\
  11 & ACH-000408 & TE-5 & Esophagus/Stomach & 0.211 & 0.582 & 0.617 & 0.035 & 0.010 & A & unexplained discrepancy & \na \\
  12 & ACH-000657 & A2780 & Ovary/Fallopian Tube & 0.230 & 0.617 & 0.577 & 0.040 & 0.013 & A & unexplained discrepancy & \na \\
  \midrule
  13 & ACH-000631 & KS-1 & CNS/Brain & 0.233 & 0.582 & 0.605 & 0.022 & 0.022 & B & unexplained discrepancy & \na \\
  14 & ACH-000149 & SK-N-SH & Peripheral Nervous System & 0.930 & 0.584 & 0.602 & 0.018 & 0.013 & B & technical review first & related lines in the cohort; Cellosaurus problem flag on a relative (CVCL\_0019); small tissue group \\
  15 & ACH-000025 & CH-157MN & CNS/Brain & 0.960 & 0.575 & 0.639 & 0.064 & 0.026 & B & unexplained discrepancy & \na \\
  16 & ACH-000545 & VM-CUB1 & Bladder/Urinary Tract & 0.920 & 0.586 & 0.608 & 0.022 & 0.016 & B & unexplained discrepancy & \na \\
  17 & ACH-000520 & 59M & Ovary/Fallopian Tube & 0.910 & 0.579 & 0.611 & 0.032 & 0.012 & B & unexplained discrepancy & \na \\
  18 & ACH-000662 & COR-L23 & Lung & 0.291 & 0.594 & 0.601 & 0.007 & 0.012 & B & unexplained discrepancy & \na \\
  19 & ACH-000423 & SK-MEL-3 & Skin & 0.880 & 0.578 & 0.539 & 0.039 & 0.023 & B & technical review first & small tissue group \\
  20 & ACH-000958 & SW48 & Bowel & 0.280 & 0.584 & 0.563 & 0.021 & 0.018 & B & unexplained discrepancy & \na \\
  21 & ACH-000652 & SUIT-2 & Pancreas & 0.885 & 0.584 & 0.562 & 0.022 & 0.006 & B & unexplained discrepancy & \na \\
  \midrule
  22 & ACH-000758 & MKN74 & Esophagus/Stomach & 0.293 & 0.582 & 0.563 & 0.019 & 0.030 & C & exploratory candidate & \na \\
  23 & ACH-000046 & ACHN & Kidney & 0.861 & 0.579 & 0.586 & 0.006 & 0.013 & C & technical review first & small tissue group \\
  24 & ACH-000511 & Calu-1 & Lung & 0.896 & 0.586 & 0.626 & 0.040 & 0.017 & C & exploratory candidate & \na \\
  25 & ACH-000984 & HEC-6 & Uterus & 0.249 & 0.572 & 0.518 & 0.054 & 0.018 & C & technical review first & small tissue group \\
  26 & ACH-000863 & DBTRG-05MG & CNS/Brain & 0.311 & 0.580 & 0.632 & 0.052 & 0.011 & C & exploratory candidate & \na \\
  27 & ACH-000538 & HUTU80 & Bowel & 0.278 & 0.582 & 0.541 & 0.041 & 0.015 & C & exploratory candidate & \na \\
  28 & ACH-000292 & NCI-H841 & Lung & 0.850 & 0.589 & 0.559 & 0.030 & 0.027 & C & exploratory candidate & \na \\
  29 & ACH-000903 & FTC-133 & Thyroid & 0.338 & 0.589 & 0.613 & 0.024 & 0.013 & C & technical review first & small tissue group \\
  30 & ACH-000940 & AN3 CA & Uterus & 0.360 & 0.576 & 0.526 & 0.051 & 0.035 & C & technical review first & small tissue group; prediction unstable across CV repeats \\
  31 & ACH-000355 & NCI-H82 & Lung & 0.730 & 0.589 & 0.517 & 0.071 & 0.020 & C & technical review first & M1 and M2 disagree strongly \\
  32 & ACH-000470 & SW1463 & Bowel & 0.685 & 0.571 & 0.534 & 0.037 & 0.036 & C & technical review first & prediction unstable across CV repeats \\
  33 & ACH-000358 & NCI-H69 & Lung & 0.425 & 0.593 & 0.543 & 0.050 & 0.038 & C & technical review first & prediction unstable across CV repeats \\
  34 & ACH-000169 & RD & Soft Tissue & 0.690 & 0.589 & 0.494 & 0.095 & 0.182 & C & technical review first & small tissue group; prediction unstable across CV repeats; M1 and M2 disagree strongly \\
  35 & ACH-000544 & OE21 & Esophagus/Stomach & 0.710 & 0.567 & 0.611 & 0.044 & 0.033 & C & technical review first & prediction unstable across CV repeats \\
  36 & ACH-000464 & CAS-1 & CNS/Brain & 0.740 & 0.575 & 0.659 & 0.083 & 0.040 & C & technical review first & prediction unstable across CV repeats; M1 and M2 disagree strongly \\
  37 & ACH-001306 & 8305C & Thyroid & 0.736 & 0.582 & 0.658 & 0.075 & 0.015 & C & technical review first & small tissue group; M1 and M2 disagree strongly \\
  38 & ACH-000703 & DMS 79 & Lung & 0.507 & 0.593 & 0.579 & 0.014 & 0.035 & C & technical review first & prediction unstable across CV repeats \\
  39 & ACH-000480 & HuH-7 & Liver & 0.510 & 0.584 & 0.652 & 0.067 & 0.051 & C & technical review first & small tissue group; prediction unstable across CV repeats; M1 and M2 disagree strongly \\
  40 & ACH-000979 & DU 145\textsuperscript{\textdagger} & Prostate & 0.710 & 0.580 & 0.650 & 0.071 & 0.048 & C & technical review first & stranded RNA profile; small tissue group; prediction unstable across CV repeats; M1 and M2 disagree strongly \\
  41 & ACH-000771 & BICR 56 & Head and Neck & 0.650 & 0.599 & 0.703 & 0.103 & 0.021 & C & technical review first & M1 and M2 disagree strongly \\
  42 & ACH-000244 & DK-MG & CNS/Brain & 0.529 & 0.574 & 0.667 & 0.094 & 0.030 & C & technical review first & M1 and M2 disagree strongly \\
  43 & ACH-000415 & BICR 6 & Head and Neck & 0.714 & 0.594 & 0.683 & 0.089 & 0.034 & C & technical review first & prediction unstable across CV repeats; M1 and M2 disagree strongly \\
  44 & ACH-000665 & SK-MES-1 & Lung & 0.675 & 0.583 & 0.647 & 0.065 & 0.033 & C & technical review first & prediction unstable across CV repeats \\
  45 & ACH-000138 & CFPAC-1 & Pancreas & 0.554 & 0.588 & 0.668 & 0.079 & 0.017 & C & technical review first & M1 and M2 disagree strongly \\
  46 & ACH-000254 & SCC-15 & Head and Neck & 0.585 & 0.598 & 0.670 & 0.072 & 0.017 & C & technical review first & M1 and M2 disagree strongly \\
  47 & ACH-000148 & Hs 578T & Breast & 0.650 & 0.566 & 0.644 & 0.078 & 0.050 & C & technical review first & prediction unstable across CV repeats; M1 and M2 disagree strongly \\
  48 & ACH-000188 & SCC-25 & Head and Neck & 0.675 & 0.598 & 0.680 & 0.082 & 0.036 & C & technical review first & prediction unstable across CV repeats; M1 and M2 disagree strongly \\
  \midrule
  206 & ACH-000019 & MCF7\textsuperscript{\textdagger} & Breast & 0.520 & 0.574 & 0.534 & 0.040 & 0.016 & \na & not prioritised & \na \\
\end{longtable}
\tablenote{\textsuperscript{\textdagger}\,Context line from the motivating \textit{Moringa}
  study (DU145 and MCF-7). \enquote{Unexplained discrepancy}: unexplained discrepancy
  requiring follow-up (tier A or B without technical flags); \enquote{exploratory candidate}:
  tier~C without technical flags; \enquote{technical review first}: at least one technical
  flag. Rank 206 of 241 for MCF-7.}
\endgroup
\end{landscape}

% =====================================================================================
\section{Compound evidence records}\label{app:compounds}

The complete curated records, with every citation, DOI, verification route, reading depth,
extraction solvent, identification method, model system, radiation dose, concentration,
measured genes and reported direction, are held in \path{data/compounds/moringa_constituents.yaml}
and \path{data/compounds/compound_biology_evidence.yaml}; the curation audit trail (searches,
exclusions, identifier problems and unresolved gaps) is in
\path{docs/compound_evidence_curation.md}. Tables~\ref{tab:c1}--\ref{tab:c4} present the
evidence used for the ranking. DOIs of the primary studies are listed in the tables; only the
studies discussed in the text appear in the reference list.

% ---------------------------------------------------------------- constituent inventory
\begin{landscape}
\begingroup\scriptsize\setlength{\tabcolsep}{3pt}
\begin{longtable}{@{}L{0.11\linewidth} r l L{0.18\linewidth} c L{0.23\linewidth}
    L{0.075\linewidth} L{0.11\linewidth} L{0.13\linewidth}@{}}
  \caption[\textit{Moringa} constituent inventory]{Inventory of documented
    \textit{M.~oleifera} seed and leaf constituents: one row per primary analytical report,
    with the extraction solvent or fraction as reported and the strength of identification
    (isolation and NMR or reference standard; spectral-library match; not stated). DCM: the
    report concerns a dichloromethane (methylene chloride) extract or fraction. Benzyl
    isothiocyanate is listed because it was evaluated, although its only report concerns its
    precursor in roots. Source:
    \path{results/tables/primary_v1/compounds_v1/table_c1_constituent_inventory.csv}.}
  \label{tab:c1}\\
  \toprule
  \thead{Compound} & \thead{\#} & \thead{Part} & \thead{Solvent or fraction} & \thead{DCM}
    & \thead{Identification method} & \thead{Strength} & \thead{Research group}
    & \thead{DOI} \\
  \midrule
  \endfirsthead
  \multicolumn{9}{@{}l}{\tcont{Table~\thetable\ (continued)}}\\
  \toprule
  \thead{Compound} & \thead{\#} & \thead{Part} & \thead{Solvent or fraction} & \thead{DCM}
    & \thead{Identification method} & \thead{Strength} & \thead{Research group}
    & \thead{DOI} \\
  \midrule
  \endhead
  \midrule
  \multicolumn{9}{@{}r}{\tcont{Continued on next page}}\\
  \endfoot
  \bottomrule
  \endlastfoot
  Acetylated moringin\textsuperscript{a} & 1 & leaf & Ethanol leaf extract & no
    & Isolation; structure elucidated by NMR & isolation, NMR & Faizi / Siddiqui (Karachi)
    & \doilink{10.1021/np50111a011} \\
   & 2 & leaf & Fresh leaves extracted with water & no & Not stated in the abstract
    & not stated & Raskin (Rutgers) & \doilink{10.1016/j.phytochem.2014.03.028} \\
  \addlinespace
  $\alpha$-Tocopherol & 1 & leaf & Leaf supercritical \ce{CO2} extract & no
    & GC-MS with NIST library and Kovats retention indices, quantified against a calibration
      curve & reference standard & Kessler (2025) & \doilink{10.1111/jocd.70486} \\
   & 2 & seed & Seed oil (cold press, \textit{n}-hexane, or chloroform/methanol 50:50) & no
    & Tocopherols by HPLC ($\alpha$, $\gamma$ and $\delta$ reported) & not stated
    & Tsaknis / Lalas & \doilink{10.1021/jf9904214} \\
  \addlinespace
  Astragalin & 1 & leaf & Not stated in the abstract (full text paywalled) & no
    & HPLC and GC-MS per MeSH; criteria not verified & not stated & Bennett / Mellon (Norwich)
    & \doilink{10.1021/jf0211480} \\
   & 2 & leaf & Leaf water extract & no & Standard purchased for quantitation; identification
      against literature MS data; recorded conservatively as a library match & library match
    & Wang J (2024) & \doilink{10.1016/j.fochx.2024.101585} \\
  \addlinespace
  Benzyl isothiocyanate & 1 & other & Roots (precursor benzyl glucosinolate), not seed or leaf
    & no & HPLC and GC-MS per MeSH & not stated & Bennett / Mellon (Norwich)
    & \doilink{10.1021/jf0211480} \\
  \addlinespace
  $\beta$-Sitosterol & 1 & seed & Ethanol seed extract & no & Isolation (also
      $\beta$-sitosterol-3-\textit{O}-glucoside and an oleoyl-glucoside); method not stated
    & not stated & Guevara / Nishino & \doilink{10.1016/s1383-5718(99)00025-x} \\
   & 2 & seed & Seed oil (cold press, \textit{n}-hexane, or chloroform/methanol 50:50) & no
    & Sterol analysis; up to 50.07\% of the sterol fraction & not stated & Tsaknis / Lalas
    & \doilink{10.1021/jf9904214} \\
   & 3 & seed & Bio-guided column fractionation; solvent not stated & no & GC-MS
    & library match & Raafat (Beirut) & \doilink{10.1007/s11655-017-2758-4} \\
  \addlinespace
  Chlorogenic acid & 1 & leaf & Not stated in the abstract (full text paywalled) & no
    & HPLC and GC-MS per MeSH; 3- and 5-caffeoylquinic acids reported & not stated
    & Bennett / Mellon (Norwich) & \doilink{10.1021/jf0211480} \\
   & 2 & leaf & Leaf water extract & no & Standard purchased for quantitation; identification
      against literature MS data; recorded as a library match & library match
    & Wang J (2024) & \doilink{10.1016/j.fochx.2024.101585} \\
  \addlinespace
  \textit{cis}-Vaccenic acid & 1 & leaf & Dichloromethane leaf extract & \textbf{yes}
    & GC-MS; 10.68\% of the reported profile; criteria not verified & library match
    & Khurana (Amity) & \doilink{10.5958/0975-6892.2019.00004.2} \\
  \addlinespace
  \textit{O}-Ethyl carbamate\textsuperscript{b} & 1 & seed & Ethanol seed extract & no
    & Reported as a new compound; method not stated & not stated & Guevara / Nishino
    & \doilink{10.1016/s1383-5718(99)00025-x} \\
   & 2 & seed & 85\% ethanol seed extract, D101 resin chromatography & no
    & Isolated as compound 7; data compared with the literature & isolation, NMR & Li LZ
    & \doilink{10.3390/molecules28176426} \\
  \addlinespace
  Glucomoringin & 1 & both & Not stated in the abstract (full text paywalled) & no
    & HPLC and GC-MS per MeSH; criteria not verified & not stated & Bennett / Mellon (Norwich)
    & \doilink{10.1021/jf0211480} \\
   & 2 & both & Not stated in the abstract & no & Not stated in the abstract & not stated
    & Bennett / Mellon (Norwich) & \doilink{10.1016/j.foodchem.2010.03.073} \\
   & 3 & seed & 85\% ethanol seed extract, D101 resin chromatography & no
    & Isolated as compound 16; data compared with the literature & isolation, NMR & Li LZ
    & \doilink{10.3390/molecules28176426} \\
   & 4 & leaf & Leaf water reflux extract & no & Orbitrap MS/MS fragmentation compared with
      the literature; no authentic standard & library match & Wang J (2024)
    & \doilink{10.1016/j.fochx.2024.101585} \\
  \addlinespace
  4-Hydroxyphenyl\-acetonitrile & 1 & leaf & Dichloromethane fraction of a leaf methanolic
      extract & \textbf{yes} & GC-MS; reported as a major compound; full list and criteria
      not verified (paywalled) & library match & Arsad (USM)
    & \doilink{10.1007/s11033-021-06466-y} \\
   & 2 & seed & Ripe seed: crude 80\% ethanol extract, butanol fraction and aqueous residue
      (DCM fraction prepared, GC-MS not reported) & no & GC-MS with NIST library match
      $>80\%$; no authentic standards & library match & Arsad (USM)
    & \doilink{10.4103/pm.pm_212_17} \\
  \addlinespace
  Isoquercitrin & 1 & leaf & Not stated in the abstract (full text paywalled) & no
    & HPLC and GC-MS per MeSH; also quercetin-3-\textit{O}-(6$''$-malonyl)glucoside
    & not stated & Bennett / Mellon (Norwich) & \doilink{10.1021/jf0211480} \\
   & 2 & leaf & Leaf water extract & no & Standard purchased for quantitation; identification
      by formula and literature MS; recorded as a library match & library match
    & Wang J (2024) & \doilink{10.1016/j.fochx.2024.101585} \\
  \addlinespace
  Kaempferol & 1 & seed & Not stated in the abstract & no & LC-HRMS dereplication (tentative
      annotation, no standard) & library match & Shady (2022)
    & \doilink{10.3390/antiox11091743} \\
  \addlinespace
  Methyl 4-hydroxyphenylacetate & 1 & leaf & Dichloromethane fraction of a leaf methanolic
      extract & \textbf{yes} & GC-MS; reported as a major compound; criteria not verified
    & library match & Arsad (USM) & \doilink{10.1007/s11033-021-06466-y} \\
  \addlinespace
  Moringin & 1 & seed & Ethanol seed extract & no & Isolation; method not stated
    & not stated & Guevara / Nishino & \doilink{10.1016/s1383-5718(99)00025-x} \\
   & 2 & leaf & Fresh leaves extracted with water & no & Not stated (full text not
      retrieved) & not stated & Raskin (Rutgers) & \doilink{10.1016/j.phytochem.2014.03.028} \\
   & 3 & leaf & Leaf fast centrifugal partition chromatography fractions & no
    & Not stated in the abstract & not stated & Raskin (Rutgers) & \doilink{10.1021/jf505014n} \\
  \addlinespace
  Niazimicin & 1 & both & \ce{CH2Cl2} partition fraction of a 70\% ethanol maceration of
      seeds and of leaves & \textbf{yes} & Reference standard isolated from the seed
      \ce{CH2Cl2} fraction (purity 96.6\% by PDA); identity by HRMS, UV, IR, EI-MS and 1D/2D
      NMR; quantified by validated UPLC-MS/MS (MRM 358\,$\rightarrow$\,107)
    & isolation, NMR & Hanafi / Abdel Motaal (Cairo) & \doilink{10.38212/2224-6614.3282} \\
   & 2 & seed & Methylene-chloride and ethyl-acetate fractions of a seed extract
    & \textbf{yes} & Not stated in the abstract & not stated & Hanafi / Abdel Motaal (Cairo)
    & \doilink{10.1111/jfbc.13992} \\
   & 3 & seed & Ethanol seed extract & no & Isolation; method not stated & not stated
    & Guevara / Nishino & \doilink{10.1016/s1383-5718(99)00025-x} \\
   & 4 & seed & Ethyl-acetate fraction of an ethanolic seed extract & no
    & Chromatographic isolation & not stated & Albuquerque Costa (Ceará)
    & \doilink{10.1155/2017/7963747} \\
   & 5 & seed & 85\% ethanol seed extract, D101 resin chromatography & no
    & Isolated as compound 11; spectroscopic data compared with the literature
    & isolation, NMR & Li LZ & \doilink{10.3390/molecules28176426} \\
  \addlinespace
  Niaziminin & 1 & leaf & Ethanol leaf extract & no & Isolation; structure elucidated by NMR
    & isolation, NMR & Faizi / Siddiqui (Karachi) & \doilink{10.1021/np50111a011} \\
   & 2 & leaf & Leaf; solvent not given in the abstract & no & Not stated in the abstract
    & not stated & Ohigashi (Kyoto) & \doilink{10.1055/s-2006-957442} \\
  \addlinespace
  Niazirin & 1 & leaf & Ethanol leaf extract & no & Bioassay-guided isolation; spectroscopy
      including 2D NMR & isolation, NMR & Faizi / Siddiqui (Karachi)
    & \doilink{10.1021/np50111a011} \\
   & 2 & seed & Ethanol seed extract & no & Isolation; method not stated & not stated
    & Guevara / Nishino & \doilink{10.1016/s1383-5718(99)00025-x} \\
   & 3 & seed & 85\% ethanol seed extract, D101 resin chromatography & no
    & Isolated as compound 9; data compared with the literature & isolation, NMR & Li LZ
    & \doilink{10.3390/molecules28176426} \\
  \addlinespace
  Oleic acid & 1 & seed & Seed oil (cold press, \textit{n}-hexane, or chloroform/methanol
      50:50) & no & Fatty-acid methyl-ester analysis; up to 75.39\% of the oil & not stated
    & Tsaknis / Lalas & \doilink{10.1021/jf9904214} \\
  \addlinespace
  Palmitic acid & 1 & leaf & Dichloromethane leaf extract & \textbf{yes} & GC-MS; 96
      compounds reported, \textit{n}-hexadecanoic acid the largest peak (28.35\%); library
      matching assumed but not verified & library match & Khurana (Amity)
    & \doilink{10.5958/0975-6892.2019.00004.2} \\
   & 2 & seed & Seed oil (cold press, \textit{n}-hexane, or chloroform/methanol 50:50) & no
    & Fatty-acid methyl-ester analysis & not stated & Tsaknis / Lalas
    & \doilink{10.1021/jf9904214} \\
  \addlinespace
  Phytol & 1 & leaf & Hydroethanolic leaf extract & no & GC-MS; 23.54\% of the leaf profile;
      method details not verified & library match & Vats / Gupta (Rajasthan)
    & \doilink{10.1007/s12298-016-0407-6} \\
  \addlinespace
  Quercetin & 1 & seed & Not stated in the abstract & no & LC-HRMS dereplication (tentative
      annotation, no standard) & library match & Shady (2022)
    & \doilink{10.3390/antiox11091743} \\
   & 2 & leaf & Leaf water extract & no & Listed in the constituent table; no authentic
      standard & library match & Wang J (2024) & \doilink{10.1016/j.fochx.2024.101585} \\
  \addlinespace
  Rutin & 1 & leaf & \textit{M.~oleifera} leaf water extract & no & Standard purchased for
      quantitation; identification against literature MS data; recorded as a library match
    & library match & Wang J (2024) & \doilink{10.1016/j.fochx.2024.101585} \\
   & 2 & leaf & Not stated in the abstract & no & Quercetin rutinosides among the leaf
      flavonoids; criteria not stated & not stated & Bennett / Mellon (Norwich)
    & \doilink{10.1016/j.foodchem.2010.03.073} \\
\end{longtable}
\tablenote{\textsuperscript{a}\,4-[(4$'$-\textit{O}-Acetyl-$\alpha$-L-rhamnosyloxy)benzyl]
  isothiocyanate. \textsuperscript{b}\,\textit{O}-Ethyl-4-($\alpha$-L-rhamnosyloxy)benzyl
  carbamate. Research groups are those used to count independent sources; reports from one
  group count once (for example the two niazimicin reports of the Cairo group, and the two
  4-hydroxyphenylacetonitrile reports of the USM group). Identity caveats recorded in the
  source file include: the glucomoringin identifier is ambiguous (every PubChem record is a
  potassium salt with a different stereo-InChIKey); free quercetin and kaempferol may arise by
  hydrolysis of glycosides; the C18:1 isomer of the \textit{cis}-vaccenic acid report cannot
  be settled by library matching; and niazirin was explicitly \emph{not} identified in the
  dichloromethane fractions analysed by \textcite{abdelsayed2021jfda}.}
\endgroup
\end{landscape}

% ------------------------------------------------------------ compound-biology evidence
\begin{landscape}
\begingroup\scriptsize\setlength{\tabcolsep}{3pt}
\begin{longtable}{@{}l L{0.075\linewidth} l L{0.065\linewidth} L{0.16\linewidth}
    L{0.085\linewidth} L{0.085\linewidth} L{0.15\linewidth} L{0.08\linewidth} c c
    L{0.11\linewidth}@{}}
  \caption[Compound--biology evidence]{Primary evidence linking the constituents to
    radiation-response biology. Level L3: ionising-radiation study of the purified compound;
    L2: purified compound modulating a relevant process without radiation. Reported directions
    are never merged: protection of normal cells, protection of tumour cells, sensitisation of
    tumour cells, process modulation without radiation, genotoxicity or no effect. Liability:
    evidence that caps the tier at C. Scored: the record counts towards the score (purified
    compound). Source:
    \path{results/tables/primary_v1/compounds_v1/table_c2_biology_evidence.csv}.}
  \label{tab:c2}\\
  \toprule
  \thead{ID} & \thead{Compound} & \thead{Level} & \thead{Process} & \thead{Model system}
    & \thead{Radiation} & \thead{Dose (Gy)} & \thead{Concentration} & \thead{Direction}
    & \thead{Liab.} & \thead{Scored} & \thead{DOI} \\
  \midrule
  \endfirsthead
  \multicolumn{12}{@{}l}{\tcont{Table~\thetable\ (continued)}}\\
  \toprule
  \thead{ID} & \thead{Compound} & \thead{Level} & \thead{Process} & \thead{Model system}
    & \thead{Radiation} & \thead{Dose (Gy)} & \thead{Concentration} & \thead{Direction}
    & \thead{Liab.} & \thead{Scored} & \thead{DOI} \\
  \midrule
  \endhead
  \midrule
  \multicolumn{12}{@{}r}{\tcont{Continued on next page}}\\
  \endfoot
  \bottomrule
  \endlastfoot
  TOC-01 & $\alpha$-Tocopherol & L3 & proliferation & Rats (Sprague-Dawley), exteriorised
      mid-small bowel & X-rays & 11 & not stated (luminal phosphate and acetate esters; also
      dietary) & protection, normal & no & yes & \doilink{10.1097/00000658-199522240-00008} \\
  TOC-02 & $\alpha$-Tocopherol & L3 & other & Humans, $n=54$, oral cavity and oropharyngeal
      cancer patients & External-beam radiotherapy & 2 per fraction over 5--7 weeks; total
      not stated & not stated (oil-solution rinse) & protection, normal & no & yes
    & \doilink{10.1002/hed.10382} \\
  AST-01 & Astragalin & L2 & DNA damage & \textit{In vitro} DNA and \textit{Saccharomyces
      cerevisiae} & none & not stated & not stated & genotoxic & \textbf{yes} & yes
    & \doilink{10.3390/ijms13032846} \\
  BIT-01 & Benzyl isothio\-cyanate & L3 & cell cycle & BxPC-3 human pancreatic cancer & Gamma
    & 5 & 2.5~µM & sensitisation, tumour & no & yes & \doilink{10.2741/e55} \\
  BIT-02 & Benzyl isothio\-cyanate & L3 & apoptosis & PANC-1 and MIAPaCa-2 human pancreatic
      cancer & X-rays & not stated & 5--10~µM (PANC-1); 2.5--5~µM (MIAPaCa-2)
    & sensitisation, tumour & no & yes & \doilink{10.3892/ijmm.2011.770} \\
  BST-01 & $\beta$-Sitosterol & L3 & DNA damage & Mouse thymocytes (BALB/c) & Ionising;
      modality unverified & 6 & not stated & protection, normal & no & yes
    & \doilink{10.1002/jcb.21326} \\
  BST-02 & $\beta$-Sitosterol & L3 & oxidative stress & Adjuvant-arthritic male rats with and
      without irradiation & Gamma & 2 per week $\times$ 4 (8 total) & 40~mg/kg/day oral
    & protection, normal & no & yes & \doilink{10.1080/08923973.2025.2591423} \\
  CGA-01 & Chlorogenic acid & L3 & oxidative stress & Hepatocellular carcinoma cell lines (not
      stated) and a murine model & Ionising & not stated & not stated
    & \textbf{protection, tumour} & \textbf{yes} & yes & \doilink{10.1155/2022/4566949} \\
  CGA-02 & Chlorogenic acid & L3 & DNA damage & Mouse bone marrow & Gamma, whole body
    & 1.15 (0.05~Gy/s) & 50, 100, 200~mg/kg oral & protection, normal & no & yes
    & \doilink{10.1016/0165-7992(93)90022-n} \\
  CGA-03 & Chlorogenic acid & L3 & apoptosis & Mice, $n=70$, testis & Ionising, total body
    & 2.0 & 5--80~mg/kg i.p. & protection, normal & no & yes
    & \doilink{10.1016/j.heliyon.2022.e10798} \\
  KAE-01 & Kaempferol & L3 & cell cycle & A-549 lung cancer and HFL1 normal lung fibroblasts;
      nude-mouse xenograft & Ionising & not stated & not stated & sensitisation, tumour & no
    & yes & \doilink{10.3892/or.2015.4204} \\
  KAE-02 & Kaempferol & L3 & apoptosis & C57BL/6 mice and HUVECs & Gamma & 8.5 and 7 total
      body; \textit{in vitro} not stated & not stated & protection, normal & no & yes
    & \doilink{10.1016/j.etap.2018.04.014} \\
  MOR-01 & Moringin & L2 & apoptosis & SH-SY5Y human neuroblastoma & none & \na & 16.4~µM
    & process, no radiation & no & yes & \doilink{10.3390/ijms20081930} \\
  MOR-02 & Moringin & L2 & cell cycle & PC-3 human prostate cancer & none & \na
    & IC\textsubscript{50} 3.5~µg/mL & process, no radiation & no & yes
    & \doilink{10.3390/molecules28073214} \\
  MOR-03 & Moringin & L2 & oxidative stress & Differentiated NSC-34 motor-neuron-like cells
    & none & \na & 0.5, 5, 10~µM ($\alpha$-cyclodextrin complex) & process, no radiation & no
    & yes & \doilink{10.3390/antiox13070813} \\
  MOR-04 & Moringin & L2 & apoptosis & Carcinoma, myeloma and leukaemia lines; nude-mouse
      xenografts & none & \na & micromolar range; not stated & process, no radiation & no
    & yes & \doilink{10.1016/j.bcp.2009.12.008} \\
  NZM-01 & Niazimicin & L2 & oxidative stress & Rats, aluminium-chloride dementia model
    & none & \na & 250~mg/kg/day (extract fraction) & process, no radiation & no
    & \textbf{no} & \doilink{10.1111/jfbc.13992} \\
  NZR-01 & Niazirin & L2 & oxidative stress & High-glucose vascular smooth muscle cells;
      diabetic mice & none & \na & not stated & process, no radiation & no & yes
    & \doilink{10.1016/j.phymed.2019.153066} \\
  OLE-01 & Oleic acid & L3 & proliferation & Human LDV lymphoid cells & Ionising (air or
      nitrogen) & not stated & not stated & \textbf{no effect} & no & yes
    & \doilink{10.1080/09553008314550431} \\
  PHY-01 & Phytol & L2 & apoptosis & A549 human lung carcinoma & none & \na & not stated
    & process, no radiation & no & yes & \doilink{10.1016/j.biopha.2017.05.066} \\
  PHY-02 & Phytol & L2 & oxidative phosphorylation & A549 human lung carcinoma & none & \na
    & not stated & process, no radiation & no & yes & \doilink{10.1016/j.biopha.2018.06.035} \\
  QUE-01 & Quercetin & L3 & DNA repair & Human ovarian cancer cells (not stated) and xenograft
    & X-rays & not stated & not stated & sensitisation, tumour & no & yes
    & \doilink{10.2147/ott.s147316} \\
  QUE-02 & Quercetin & L3 & proliferation & Medulloblastoma lines versus neural precursors and
      normal fibroblasts; orthotopic xenografts & Ionising, fractionated & not stated & low
      micromolar & sensitisation, tumour & no & yes & \doilink{10.18632/oncotarget.7980} \\
  QUE-03 & Quercetin & L3 & DNA damage & Human peripheral blood lymphocytes; plasmid DNA
    & Gamma & 1--4 (cells); 6 (plasmid) & 3--48~µM; 24~µM optimal & protection, normal & no
    & yes & \doilink{10.1016/j.mrgentox.2008.03.003} \\
  QUE-04 & Quercetin & L3 & other & C3H/HeN mouse skin, human and murine fibroblasts, murine
      tumours & X-rays & 35 & 1\% in chow & protection, normal & no & yes
    & \doilink{10.1667/rr3237.1} \\
  RUT-01 & Rutin & L3 & DNA damage & Human peripheral blood lymphocytes & Gamma & 2
    & 25~µg/mL & protection, normal & no & yes & \doilink{10.4103/ijnm.ijnm_30_17} \\
  RUT-02 & Rutin & L3 & oxidative stress & Rat brain & Gamma & 5 & 200~mg/kg/day oral
    & protection, normal & no & yes & \doilink{10.1080/13813455.2017.1374978} \\
\end{longtable}
\tablenote{NZM-01 is recorded so that the absence of purified-niazimicin biology is explicit:
  the dosed material was an extract fraction, not the purified compound, so the record does
  not score. CGA-01 is liability evidence and is never pooled with the normal-tissue
  chlorogenic-acid studies. QUE-04 reports no effect on fibroblast or murine tumour
  radioresponse. Several abstracts omit cell lines, doses or concentrations; these gaps are
  recorded as \enquote{not stated} and must be recovered from the full texts.}
\endgroup
\end{landscape}

% ---------------------------------------------------------- ranking justifications
\begin{landscape}
\begingroup\scriptsize\setlength{\tabcolsep}{3pt}
\begin{longtable}{@{}L{0.11\linewidth} L{0.15\linewidth} L{0.19\linewidth}
    L{0.27\linewidth} L{0.11\linewidth} L{0.15\linewidth}@{}}
  \caption[Written justification of every ranking score]{Written justification of every
    point of the candidate ranking (Table~\ref{tab:ranking}) and of each tier, as recorded by
    the scoring code. Source:
    \path{results/tables/primary_v1/compounds_v1/table_c3_candidate_ranking.csv}.}
  \label{tab:c3reasons}\\
  \toprule
  \thead{Compound} & \thead{$D_1$ identity} & \thead{$D_2$ radiation biology}
    & \thead{$D_3$ independent sources} & \thead{$D_4$ feasibility} & \thead{Tier reason} \\
  \midrule
  \endfirsthead
  \multicolumn{6}{@{}l}{\tcont{Table~\thetable\ (continued)}}\\
  \toprule
  \thead{Compound} & \thead{$D_1$ identity} & \thead{$D_2$ radiation biology}
    & \thead{$D_3$ independent sources} & \thead{$D_4$ feasibility} & \thead{Tier reason} \\
  \midrule
  \endhead
  \midrule
  \multicolumn{6}{@{}r}{\tcont{Continued on next page}}\\
  \endfoot
  \bottomrule
  \endlastfoot
  $\alpha$-Tocopherol & 1 primary report with structural confirmation (reference standard)
    & Ionising-radiation evidence from 2 independent groups (Delaney, Minnesota; Ferreira,
      Brazil) & 2 groups reporting it in \textit{Moringa}, 2 in the biology evidence, 4
      distinct in total & Identifier verified; 2 scored studies report a usable
      concentration & Total 8/8 with $D_1\ge1$ and $D_2\ge2$ \\
  Chlorogenic acid & 2 primary reports; identification by spectral-library match or not
      stated & Ionising-radiation evidence from 3 independent groups (Abraham / Kesavan; Shoorei
      / Karimipour; Yan) & 2 groups in \textit{Moringa}, 3 in the biology evidence, 5 distinct
    & Identifier verified; 3 scored studies report a usable concentration
    & Capped at C from A (total 7/8): liability evidence CGA-01, protection of tumour cells \\
  Kaempferol & 1 primary report; spectral-library match or not stated & Ionising-radiation
      evidence from 2 independent groups (Yao, Taiwan; Yuan, China) & 1 group in
      \textit{Moringa}, 2 in the biology evidence, 3 distinct & Identifier verified; 2 scored
      studies report a usable concentration & Total 7/8 with $D_1\ge1$ and $D_2\ge2$ \\
  Quercetin & 2 primary reports; spectral-library match or not stated & Ionising-radiation
      evidence from 4 independent groups (Citrin, NCI; Gong / Liu, China; Hulleman /
      Wurdinger, Amsterdam; Menon, Annamalai) & 2 groups in \textit{Moringa}, 4 in the
      biology evidence, 6 distinct & Identifier verified; 4 scored studies report a usable
      concentration & Total 7/8 with $D_1\ge1$ and $D_2\ge2$ \\
  Rutin & 2 primary reports; spectral-library match or not stated & Ionising-radiation
      evidence from 2 independent groups (Abdel-Rafei / Thabet, NCRRT Egypt; Patil,
      Mangalore) & 2 groups in \textit{Moringa}, 2 in the biology evidence, 4 distinct
    & Identifier verified; 2 scored studies report a usable concentration
    & Total 7/8 with $D_1\ge1$ and $D_2\ge2$ \\
  $\beta$-Sitosterol & 3 primary reports; spectral-library match or not stated
    & Ionising-radiation evidence from 2 independent groups (Abdel-Rafei / Thabet, NCRRT
      Egypt; Wang SQ, China) & 3 groups in \textit{Moringa}, 2 in the biology evidence, 5
      distinct & Identifier verified; 2 scored studies report a usable concentration
    & Total 7/8 with $D_1\ge1$ and $D_2\ge2$ \\
  Benzyl isothiocyanate & No primary analytical study identifies it in \textit{Moringa}
      seed or leaf (review-only, tentative or another plant part) & Ionising-radiation
      evidence from 2 independent groups (Kubota / Okayasu; Srivastava) & 1 group in
      \textit{Moringa}, 2 in the biology evidence, 3 distinct & Identifier verified; 2
      scored studies report a usable concentration & Not assessable: no primary
      \textit{Moringa} identification (prespecified rule) \\
  Niazirin & 2 primary reports with structural confirmation (isolation, NMR) & No
      ionising-radiation study; 1 study of a relevant process without radiation (Weng, China)
    & 3 groups in \textit{Moringa}, 1 in the biology evidence, 4 distinct & Identifier
      verified; 1 scored study reports a usable concentration & Total 6/8 with $D_1\ge1$ and
      $D_2\ge1$ \\
  Oleic acid & 1 primary report; spectral-library match or not stated & Ionising-radiation
      evidence from one group (Cramp, UK) & 1 group in \textit{Moringa}, 1 in the biology
      evidence, 2 distinct & Identifier verified; 1 scored study reports a usable
      concentration & Total 5/8 with $D_1\ge1$ and $D_2\ge1$ \\
  Astragalin & 2 primary reports; spectral-library match or not stated & No
      ionising-radiation study; 1 study of a relevant process without radiation (Saito)
    & 2 groups in \textit{Moringa}, 1 in the biology evidence, 3 distinct & Identifier
      verified; 1 scored study reports a usable concentration & Capped at C from B (total
      5/8): liability evidence AST-01, genotoxic \\
  Moringin & 3 primary reports; spectral-library match or not stated & No ionising-radiation
      study; 4 studies of relevant processes without radiation (Abdull Razis; D'Incalci;
      Mazzon; Navarra) & 2 groups in \textit{Moringa}, 4 in the biology evidence, 6 distinct
    & Identifier verified; 4 scored studies report a usable concentration & Total 5/8 with
      $D_1\ge1$ and $D_2\ge1$ \\
  Phytol & 1 primary report; spectral-library match or not stated & No ionising-radiation
      study; 2 studies of relevant processes without radiation (Devi; Thakkar) & 1 group in
      \textit{Moringa}, 2 in the biology evidence, 3 distinct & Identifier verified; 2 scored
      studies report a usable concentration & Total 5/8 with $D_1\ge1$ and $D_2\ge1$ \\
  Glucomoringin & 1 primary report with structural confirmation (isolation, NMR) & No primary
      study of the purified compound in a relevant process & 3 groups in \textit{Moringa}, 0
      in the biology evidence, 3 distinct & No concentration reported in any scored study
    & Total 4/8 with $D_1\ge1$ and $D_2\ge0$ \\
  Niazimicin & 2 primary reports with structural confirmation (isolation, NMR) & No primary
      study of the purified compound in a relevant process & 4 groups in \textit{Moringa}, 0
      in the biology evidence, 4 distinct & No concentration reported in any scored study
    & Total 4/8 with $D_1\ge1$ and $D_2\ge0$ \\
  Acetylated moringin & 1 primary report with structural confirmation (isolation, NMR) & No
      primary study of the purified compound in a relevant process & 2 groups in
      \textit{Moringa}, 0 in the biology evidence, 2 distinct & No concentration reported in
      any scored study & Total 3/8 with $D_1\ge1$ and $D_2\ge0$ \\
  Niaziminin & 1 primary report with structural confirmation (isolation, NMR) & No primary
      study of the purified compound in a relevant process & 2 groups in \textit{Moringa}, 0
      in the biology evidence, 2 distinct & No concentration reported in any scored study
    & Total 3/8 with $D_1\ge1$ and $D_2\ge0$ \\
  \textit{O}-Ethyl carbamate & 1 primary report with structural confirmation (isolation, NMR)
    & No primary study of the purified compound in a relevant process & 2 groups in
      \textit{Moringa}, 0 in the biology evidence, 2 distinct & No concentration reported in
      any scored study & Total 3/8 with $D_1\ge1$ and $D_2\ge0$ \\
  Isoquercitrin & 2 primary reports; spectral-library match or not stated & No primary study
      of the purified compound in a relevant process & 2 groups in \textit{Moringa}, 0 in the
      biology evidence, 2 distinct & No concentration reported in any scored study & Total
      2/8 with $D_1\ge1$ and $D_2\ge0$ \\
  Palmitic acid & 2 primary reports; spectral-library match or not stated & No primary study
      of the purified compound in a relevant process & 2 groups in \textit{Moringa}, 0 in the
      biology evidence, 2 distinct & No concentration reported in any scored study & Total
      2/8 with $D_1\ge1$ and $D_2\ge0$ \\
  4-Hydroxyphenyl\-acetonitrile & 2 primary reports; spectral-library match or not stated
    & No primary study of the purified compound in a relevant process & At most one
      independent group (Arsad, USM) & No concentration reported in any scored study
    & Total 1/8, below the tier-C threshold \\
  Methyl 4-hydroxyphenyl\-acetate & 1 primary report; spectral-library match or not stated
    & No primary study of the purified compound in a relevant process & At most one
      independent group (Arsad, USM) & No concentration reported in any scored study
    & Total 1/8, below the tier-C threshold \\
  \textit{cis}-Vaccenic acid & 1 primary report; spectral-library match or not stated & No
      primary study of the purified compound in a relevant process & At most one independent
      group (Khurana, Amity) & No concentration reported in any scored study & Total 1/8,
      below the tier-C threshold \\
\end{longtable}
\endgroup
\end{landscape}

% -------------------------------------------------------------- pathway overlap
\begingroup\footnotesize
\begin{longtable}{@{}l l S[table-format=1.0] S[table-format=2.0] S[table-format=3.0]
    L{0.21\linewidth} L{0.13\linewidth}@{}}
  \caption[Compound--pathway measured-gene overlap]{Overlap between the genes a published
    experiment measured for a compound and the members of a tier-C Hallmark gene set (prefix
    \texttt{HALLMARK\_} omitted). An overlap marks where a compound and a model-highlighted
    gene set touch the same genes, which is a hypothesis to test; it is not evidence that the
    compound modifies the gene set, the pathway or clonogenic SF2, and the gene-set ranking
    itself is capped at tier~C. Benzyl isothiocyanate is not assessable and is not in the
    shortlist. Source:
    \path{results/tables/primary_v1/compounds_v1/table_c4_pathway_overlap.csv}.}
  \label{tab:c4}\\
  \toprule
  \thead{Compound} & \thead{Gene set} & {\thead{Overlap}} & {\thead{Measured}}
    & {\thead{In set}} & \thead{Overlapping genes} & \thead{Evidence} \\
  \midrule
  \endfirsthead
  \multicolumn{7}{@{}l}{\tcont{Table~\thetable\ (continued)}}\\
  \toprule
  \thead{Compound} & \thead{Gene set} & {\thead{Overlap}} & {\thead{Measured}}
    & {\thead{In set}} & \thead{Overlapping genes} & \thead{Evidence} \\
  \midrule
  \endhead
  \midrule
  \multicolumn{7}{@{}r}{\tcont{Continued on next page}}\\
  \endfoot
  \bottomrule
  \endlastfoot
  Benzyl isothiocyanate & P53\_PATHWAY & 2 & 9 & 200 & APAF1; CDKN1A & BIT-01; BIT-02 \\
   & UV\_RESPONSE\_DN & 2 & 9 & 144 & MAPK14; NFKB1 & BIT-01 \\
  \addlinespace
  $\beta$-Sitosterol & OXIDATIVE\_PHOSPHORYLATION & 2 & 10 & 200 & AIFM1; CYCS & BST-01 \\
   & P53\_PATHWAY & 2 & 10 & 200 & FOS; JUN & BST-01 \\
   & UV\_RESPONSE\_UP & 2 & 10 & 158 & CASP3; FOS & BST-01 \\
   & COMPLEMENT & 1 & 10 & 200 & CASP3 & BST-01 \\
  \addlinespace
  Chlorogenic acid & OXIDATIVE\_PHOSPHORYLATION & 1 & 3 & 200 & BAX & CGA-03 \\
   & P53\_PATHWAY & 1 & 3 & 200 & BAX & CGA-03 \\
  \addlinespace
  Kaempferol & COMPLEMENT & 2 & 4 & 200 & CASP3; CASP9 & KAE-02 \\
   & OXIDATIVE\_PHOSPHORYLATION & 1 & 4 & 200 & CYCS & KAE-02 \\
   & UV\_RESPONSE\_UP & 1 & 4 & 158 & CASP3 & KAE-02 \\
  \addlinespace
  Moringin & P53\_PATHWAY & 3 & 11 & 200 & BAX; CDKN1A; TP53 & MOR-01; MOR-02; MOR-03 \\
   & COMPLEMENT & 2 & 11 & 200 & CASP3; CASP9 & MOR-01; MOR-03 \\
   & OXIDATIVE\_PHOSPHORYLATION & 1 & 11 & 200 & BAX & MOR-01; MOR-02; MOR-03 \\
   & UV\_RESPONSE\_DN & 1 & 11 & 144 & NFKB1 & MOR-01; MOR-04 \\
   & UV\_RESPONSE\_UP & 1 & 11 & 158 & CASP3 & MOR-01; MOR-03 \\
  \addlinespace
  Phytol & COMPLEMENT & 2 & 5 & 200 & CASP3; CASP9 & PHY-01 \\
   & P53\_PATHWAY & 1 & 5 & 200 & FAS & PHY-01 \\
   & UV\_RESPONSE\_UP & 1 & 5 & 158 & CASP3 & PHY-01 \\
  \addlinespace
  Quercetin & P53\_PATHWAY & 5 & 9 & 200 & BAX; CDKN1A; DDIT3; TGFB1; TP53 & QUE-01; QUE-04 \\
   & OXIDATIVE\_PHOSPHORYLATION & 1 & 9 & 200 & BAX & QUE-01 \\
   & UV\_RESPONSE\_UP & 1 & 9 & 158 & H2AX & QUE-01 \\
  \addlinespace
  Rutin & KRAS\_SIGNALING\_UP & 2 & 7 & 200 & IL1B; NGF & RUT-02 \\
   & COMPLEMENT & 1 & 7 & 200 & IL6 & RUT-02 \\
   & UV\_RESPONSE\_UP & 1 & 7 & 158 & IL6 & RUT-02 \\
\end{longtable}
\endgroup

% =====================================================================================
\section{Data provenance and table correspondence}\label{app:provenance}

Dataset manifests are held in \path{data/manifests/}, one file per source and release, with
DOI, URL, licence, access date, per-file SHA-256, column meanings, units, endpoint definition
and transformations. The run records are \path{results/models/primary_v1/run_manifest.json},
\path{results/tables/primary_v1/evaluation_record.json},
\path{results/triage/primary_v1/triage_manifest.json} and the compound record
\path{results/tables/primary_v1/compounds_v1/compound_record.json}; each records the git
commit, a flag for uncommitted changes, configuration and file checksums, package versions
and UTC timestamps. Table~\ref{tab:provenance} maps every result table of this report to the
project's table identifier and source file.

\begin{table}[htbp]
  \caption{Correspondence between the result tables of this report, the project's table
    identifiers and the saved output files.}\label{tab:provenance}
  \footnotesize
  \begin{tabular}{@{}l l L{0.66\linewidth}@{}}
    \toprule
    \thead{This report} & \thead{Project ID} & \thead{Source file} \\
    \midrule
    Table~\ref{tab:composition} & Table 1 & \path{results/tables/primary_v1/publication_v1/table_1_dataset_composition.csv} \\
    Table~\ref{tab:performance} & Table 2 & \path{results/tables/primary_v1/publication_v1/table_2_model_comparison.csv} \\
    Table~\ref{tab:tissue} & Table 3 & \path{results/tables/primary_v1/publication_v1/table_3_tissue_performance.csv} \\
    Table~\ref{tab:r1} & Table R1 & \path{results/tables/exploratory_m3_v1/model_metrics.csv} \\
    Table~\ref{tab:r2} & Table R2 & \path{results/tables/loto_v1/held_out_tissue.csv} \\
    Table~\ref{tab:r3} & Table R3 & \path{results/tables/scales_v1/comparisons.csv} \\
    Table~\ref{tab:ranking} & Table C3 & \path{results/tables/primary_v1/compounds_v1/table_c3_candidate_ranking.csv} \\
    Table~\ref{tab:s1} & Table S1 & \path{results/tables/primary_v1/publication_v1/table_s1_repeat_delta_mae.csv} \\
    Table~\ref{tab:s4} & Table S4 & \path{results/tables/primary_v1/publication_v1/within_tissue_spearman_v2.csv} \\
    Table~\ref{tab:s2} & Table S2 & \path{results/tables/primary_v1/publication_triage_v1/table_s2_gene_set_ranking.csv} \\
    Table~\ref{tab:s3} & Table S3 & \path{results/tables/primary_v1/publication_triage_v1/table_s3_discrepancy_review.csv} \\
    Table~\ref{tab:c1} & Table C1 & \path{results/tables/primary_v1/compounds_v1/table_c1_constituent_inventory.csv} \\
    Table~\ref{tab:c2} & Table C2 & \path{results/tables/primary_v1/compounds_v1/table_c2_biology_evidence.csv} \\
    Table~\ref{tab:c3reasons} & Table C3 & \path{results/tables/primary_v1/compounds_v1/table_c3_candidate_ranking.csv} (justification columns) \\
    Table~\ref{tab:c4} & Table C4 & \path{results/tables/primary_v1/compounds_v1/table_c4_pathway_overlap.csv} \\
    \bottomrule
  \end{tabular}
  \tablenote{Numbers are rounded for display; the source files hold full precision. A small
    number of columns of the source files (identifiers such as PubChem CIDs and InChIKeys,
    percentile columns and verbatim quotations) are not reproduced here.}
\end{table}

% =====================================================================================
\section{Configuration, commands and documentation}\label{app:config}

\begin{table}[htbp]
  \caption{Configuration file and script for each step of the pipeline.}\label{tab:config}
  \footnotesize
  \begin{tabular}{@{}L{0.30\linewidth}L{0.30\linewidth}L{0.34\linewidth}@{}}
    \toprule
    \thead{Step} & \thead{Configuration} & \thead{Script} \\
    \midrule
    Preprocessing and quality control & \path{configs/preprocessing.yaml} & \path{scripts/run_preprocessing.py} \\
    Primary experiment & \path{configs/experiment.yaml} & \path{scripts/run_model.py} \\
    Evaluation & \path{configs/experiment.yaml} & \path{scripts/evaluate_models.py} \\
    Explainability & \path{configs/experiment.yaml} & \path{scripts/run_explainability.py} \\
    Pathway and cell-line triage & \path{configs/triage.yaml} & \path{scripts/generate_triage.py} \\
    Publication tables & \path{configs/publication.yaml} & \path{scripts/make_publication_tables.py} \\
    Exploratory M3 & \path{configs/experiment_exploratory_m3.yaml} & \path{scripts/run_model.py} \\
    Leave-one-tissue-out & \path{configs/experiment_loto.yaml} & \path{scripts/run_model.py} \\
    Outcome scales & \path{configs/experiment_scales.yaml} & \path{scripts/run_model.py} \\
    Compound evidence and ranking & \path{configs/compounds.yaml} & \path{scripts/make_compound_tables.py} \\
    Report figures & \path{configs/figures_report.yaml} & \path{scripts/make_report_figures.py} \\
    Negative controls (not executed) & \path{configs/experiment_controls.yaml} & \path{scripts/run_model.py} \\
    Coverage simulation (deferred) & \path{configs/simulation.yaml} & \path{scripts/run_interval_simulation.py} \\
    \bottomrule
  \end{tabular}
\end{table}

Methodological documentation: study protocol \path{docs/protocol.md}; architecture and naming
\path{docs/architecture.md}; data rules \path{docs/data_management.md}; figure standards
\path{docs/figure_standards.md}; compound prioritisation rules \path{docs/compound_protocol.md};
robustness plan \path{docs/robustness_plan.md}; reviewer matrix
\path{docs/agent_workflow.md}; phase gates \path{docs/research_roadmap.md}; and every decision,
open question and investigator ruling in \path{docs/decision_log.md}. The page-by-page
extraction of the original manuscript is \path{docs/moringa_scientific_linkage.md}.

% =====================================================================================
\section{Reproducibility information}\label{app:reproducibility}

The environment is specified in \path{environment.yml} (conda environment
\path{explainable-radiosensitivity-ml}, Python~3.11); package versions are recorded in every
run record but are not yet pinned in a lockfile. The primary run used six physical CPU cores
(twelve logical threads) and about 31~GiB of memory, with bounded multiprocessing and one BLAS
thread per worker; the base random seed is 20261021. Tests are in \path{tests/}; tests that
would fit a model are marked \path{model_fit} and excluded from routine runs. Result folders
and figures are write-once, so regeneration requires an explicit version change or an explicit
overwrite flag.


\end{document}
